\chapter{Tiên đề của lí thuyết tập hợp}

\

Bạn đọc đã có thể từng học qua lí thuyết tập hợp từ cấp học phổ thông và được thực hiện xây dựng các tập hợp như
\begin{itemize}
   \item nhóm các nhân vật trong Đô-rê-mon: \{Nô-bi-ta; Xu-ka; Xê-kô; Chai-en; Đô-rê-mon; \dots\}, hay
   \item tập hợp các số tự nhiên: $\{0; 1; 2; 3; 4; \dots\}$.
\end{itemize}
Một định nghĩa tương đối ngây thơ có thể được đưa ra, cho phép một tập hợp có thể là một nhóm của các sự vật, hiện tượng cụ thể hay trừu tượng, yêu cầu duy nhất chúng ta cần thỏa mãn là có thể xác định được một cách rõ ràng một đối tượng có thuộc tập hợp đó hay không. Đây gọi là \defText{nguyên lí bao hàm không hạn chế}, viết dưới dạng ngôn ngữ lô-gích: Một tập hợp là một nhóm các hạng thức $x$ thỏa mãn $A(x)$ với $A$ là một vị từ chỉ phụ thuộc vào hạng thức tự do $x$. Không may, cách tiếp cận này dẫn đến những mâu thuẫn nội tại, nổi tiếng nhất là \defText{nghịch lí Rất-xen}\footnote{Bertrand Arthur William Russell (1872 -- 1970).}. Xét vị từ ``$x$ không thuộc vào tập hợp $x$''. Theo nguyên lí bao hàm, tồn tại một tập hợp $B$ bao gồm tất cả các tập hợp $x$ thỏa mãn $x \notin x$. Vậy, $B$ có thuộc $B$ hay không?
\begin{itemize}
   \item Nếu có, thì theo đúng điều kiện xác định của tập con, $B$ không thể tự chứa nó, suy ra $B \notin B$. Đây là một mâu thuẫn.
   \item Nếu không, thì vì $B$ không tự chứa chính nó, nó hoàn toàn thỏa mãn vị từ để trở thành một phần tử thuộc về tập hợp $B$, dẫn đến việc $B \in B$. Đây cũng là một mâu thuẫn!
\end{itemize}
Tình huống nan giải này khiến định nghĩa ngây thơ ban đầu bị phá sản, buộc các nhà toán học phải thiết lập một hệ tiên đề chặt chẽ và cẩn thận hơn nhằm ngăn chặn những mâu thuẫn như trên. Sau một quá trình phát triển, hệ tiên đề ZFC\footnote{Txéc-mê-lô -- Phren-ken (đặt tên theo Ernst Friedrich Ferdinand Zermelo (1871 -- 1953) và Abraham Fraenkel (1891 -- 1965)) cùng với tiên đề chọn (axiom of Choice).} đạt được nhiều thành công, và cũng sẽ là nền tảng mà chúng ta sẽ quan tâm.

Trong chương này, khi viết $\forall$ và $\exists$, thì hạng thức đi kèm theo đó mặc định là một tập hợp.

\section{Tiên đề ngoại diên}

\

Mặc dù cái tên có vẻ lạ, ý nghĩa của nó là một tính chất vô cùng quen thuộc của tập hợp: hai tập hợp cùng phần tử thì bằng nhau. Chúng ta phát biểu nó như sau:
$$
    \defMath{\forall x, \forall y, \left( x = y \iff  \forall z, \left( z \in x \iff z \in y \right) \right)}.
$$
Hai tập bằng nhau sẽ được thừa hưởng các tính chất cho bởi phép đồng nhất từ giải tích vị từ bậc nhất.

\begin{figure}[H]
    \begin{center}
        \begin{tikzpicture}[>=stealth]
            % Vẽ hai tập hợp
            \draw[thick] (-2,0) ellipse (1.2cm and 2cm);
            \node at (-2, 2.4) {Tập hợp $x$};

            \draw[thick] (2,0) ellipse (1.2cm and 2cm);
            \node at (2, 2.4) {Tập hợp $y$};

            % Các phần tử trong tập x
            \node[element, label=left:$z_1$] (x1) at (-2, 1) {};
            \node[element, label=left:$z_2$] (x2) at (-2, 0) {};
            \node[element, label=left:$z_3$] (x3) at (-2, -1) {};

            % Các phần tử trong tập y
            \node[element, label=right:$z_1$] (y1) at (2, 1) {};
            \node[element, label=right:$z_3$] (y2) at (2, 0) {};
            \node[element, label=right:$z_2$] (y3) at (2, -1) {};

            % Nối mũi tên hai chiều
            \draw[guide double arrow] (x1) to[bend left=20] (y1);
            \draw[guide double arrow] (x2) to (y3);
            \draw[guide double arrow] (x3) to[bend left=10] (y2);
        \end{tikzpicture}
    \end{center}
    \caption{Minh họa hai tập hợp bằng nhau $x = y = \{z_1; z_2; z_3\}$.}
\end{figure}

Trình bày bằng lời, chúng ta coi hai tập hợp là bằng nhau nếu mỗi phần tử trong một tập hợp thì đều thuộc tập hợp còn lại, và ngược lại. Trong trường hợp mà chúng ta chỉ có một chiều, thì chúng ta có định nghĩa \defText{tập con}:
$$
    \defMath{\forall x, \forall y, \left( x \subseteq y \iff  \forall z, \left( z \in x \implies z \in y \right) \right)}.
$$


\begin{figure}[H]
    \begin{center}
        \begin{tikzpicture}[>=stealth]
            % Vẽ tập hợp x
            \draw[thick] (-2,0) ellipse (1.2cm and 2cm);
            \node at (-2, 2.4) {Tập hợp $x$};

            % Vẽ tập hợp y nghiêng và to hơn
            \draw[thick, rotate around={-20:(2.5,0)}] (2.5,0) ellipse (1.6cm and 2.5cm);
            \node at (2.5, 2.8) {Tập hợp $y$};

            % Các phần tử trong tập x
            \node[element, label=left:$z_1$] (x1) at (-2, 1) {};
            \node[element, label=left:$z_2$] (x2) at (-2, 0) {};
            \node[element, label=left:$z_3$] (x3) at (-2, -1) {};

            % Các phần tử trong tập y (nhiều phần tử hơn, phân bố bên trong elip nghiêng)
            \node[element, label=right:$z_1$] (y1) at (1.8, 1) {};
            \node[element, label=right:$z_2$] (y2) at (1.8, -0.2) {};
            \node[element, label=right:$z_3$] (y3) at (2.5, -1.2) {};
            \node[element, label=right:$z_4$] (y4) at (3.2, 0.5) {};
            \node[element, label=right:$z_5$] (y5) at (2.8, 1.6) {};

            % Nối mũi tên một chiều từ x sang y
            \draw[guide arrow] (x1) to[bend left=15] (y1);
            \draw[guide arrow] (x2) to (y2);
            \draw[guide arrow] (x3) to[bend right=15] (y3);
        \end{tikzpicture}
    \end{center}
    \caption{Minh họa tập con $x \subseteq y$ hay $y \supseteq x$.}
\end{figure}
Ngoài định nghĩa cho phép $\subseteq$, chúng ta còn có các định nghĩa sau: Với mọi tập hợp $x$ và $y$,
\begin{itemize}
   \item $\defMath{x \subseteqq y \iff x \subseteq y}$,
   \item $\defMath{x \supseteq y \iff y \subseteq x}$,
   \item $\defMath{x \supseteqq y \iff x \supseteq y}$,
   \item $\defMath{x \subset y \iff x \subseteq y \land x \neq y}$\footnote{Trong nhiều tài liệu, $\subset$ có ý nghĩa tương đương với $\subseteq$.},
   \item $\defMath{x \subsetneq y \iff x \subset y}$,
   \item $\defMath{x \subsetneqq y \iff x \subsetneq y}$,
   \item $\defMath{x \supset y \iff x \supseteq y \land x \neq y}$,
   \item $\defMath{x \supsetneq y \iff x \supset y}$,
   \item $\defMath{x \supsetneqq y \iff x \supsetneq y}$.
\end{itemize}

Nếu $x$ không phải là tập con của $y$ thì sao? Chúng ta có biến đổi như sau: Với mọi tập hợp $x$ và $y$,
\begin{align*}
    x \not \subseteq y & \iff \overline{\forall z, \left( z \in x \implies z \in y \right)}; \\
    x \not \subseteq y & \iff \exists z, \overline{\left( z \in x \implies z \in y \right)}; \\
    x \not \subseteq y & \iff \exists z, \left( z \in x \land z \notin y \right).
\end{align*}
Tương tự với định nghĩa $\supseteq$, chúng ta cũng có $\forall x, \forall y, (x \not \supseteq y \iff y \not \subseteq x)$. Hình \ref{fig:li_thuyet_tap_hop:tien_de:khong_la_tap_con} minh họa cho $x \not \subseteq y$.

\begin{figure}[H]
   \begin{center}
      \begin{tikzpicture}[>=stealth]
         % Vẽ tập hợp x
         \draw[thick] (-2,0) ellipse (1.2cm and 2cm);
         \node at (-2, 2.4) {Tập hợp $x$};

         % Vẽ tập hợp y
         \draw[thick, rotate around={20:(2.5,0)}] (2.5,0) ellipse (1.6cm and 2.5cm);
         \node at (2.5, 2.8) {Tập hợp $y$};

         % Các phần tử trong tập x
         \node[element, label=left:$z_1$] (x1) at (-2, 1) {};
         \node[element, label=left:$z_2$] (x2) at (-2, 0) {};
         \node[element emphasis, label={[colorEmphasis]left:$z_3$}] (x3) at (-2, -1) {};

         % Các phần tử trong tập y
         \node[element, label=right:$z_1$] (y1) at (1.8, 1) {};
         \node[element, label=right:$z_4$] (y2) at (1.8, 0) {};
         \node[element, label=right:$z_2$] (y4) at (3.2, 0.5) {};
         \node[element, label=right:$z_5$] (y5) at (2.8, -1.0) {};

         % Nối mũi tên một chiều từ x sang y cho z_1, z_2
         \draw[guide arrow] (x1) to[bend left=15] (y1);
         \draw[guide arrow] (x2) to[bend left=15] (y4);

         % Làm nổi bật rằng z không có tương ứng trong y
         \node[colorEmphasis] (text) at (0.25, -1) {$z_3 \notin y$};
         \draw[guide arrow, colorEmphasis, densely dashed] (x3) -- (text);
      \end{tikzpicture}
   \end{center}
   \caption{Minh họa $x \not \subseteq y$ (tồn tại $z_3 \in x$ nhưng $z_3 \notin y$).}
   \label{fig:li_thuyet_tap_hop:tien_de:khong_la_tap_con}
\end{figure}

Trước khi đi đến với những tiên đề tiếp theo, chúng ta sẽ cùng tìm hiểu về các tính chất mới mà phép đồng nhất có được từ tiên đề ngoại diên. Thứ nhất, để ý rằng, do $A \iff A$ đúng với mọi $A$ từ giải tích vị từ cho nên, với
\begin{equation*}
    \forall z, z \in x \iff z \in x.
\end{equation*}
Do đó, chúng ta có \defText{tính phản xạ}:
\begin{equation*}
    \defMath{\forall x, x = x}.
\end{equation*}

Ngoài ra, phép $=$ cũng có \defText{tính đối xứng}. Để ý rằng, do tính giao hoán của phép đẳng giá, chúng ta có với mọi $x$ và $y$.
\begin{equation*}
    (\forall z, z \in x \iff z \in y) \iff (\forall z, z \in y \iff z \in x).
\end{equation*}
Sử dụng định nghĩa của phép đồng nhất để được
\begin{equation*}
    \defMath{\forall x, \forall y, x = y \iff y = x}.
\end{equation*}

Một tính chất cũng quan trọng khác, với mọi $w, x, y$, chúng ta có khẳng định sau theo định nghĩa:
\begin{equation*}
    x = w \land w = y \iff \begin{cases}
        \forall z, z \in x \iff z \in w \\
        \forall z, z \in w \iff z \in y
    \end{cases}.
\end{equation*}
Theo tính chất phân phối của lượng từ toàn cục,
\begin{align*}
    x = w \land w = y \iff\forall z, \big((z \in x \iff z \in w) \land (z \in w \iff z \in y)\big).
\end{align*}
Tính chất bắc cầu của phép đẳng giá cho
\begin{equation*}
    x = w \land w = y \implies \forall z, z \in x \iff z \in y.
\end{equation*}
Và qua đó,
\begin{equation*}
    \defMath{\forall w, \forall x, \forall y, \left( x = w \land w = y \implies x = y \right)}.
\end{equation*}
Đây là tính chất còn lại, \defText{tính bắc cầu}.

\exercise Cho $A$, $B$ và $C$ là các tập hợp. Chứng minh rằng
\begin{enumerate}
   \item $A \subseteq A$;
   \item $A = B \iff A \subseteq B \land B \subseteq A$;
   \item $A \neq B \iff \exists z, (z \in A \land z \notin B) \lor (z \notin A \land z \in B)$;
   \item $A \subseteq B \land B \subseteq C \implies A \subseteq C$.
\end{enumerate}

\solution

\setcounter{subexercise}{1}
\arabic{subexercise}. Có $\forall z, (z \in A \implies z \in A)$ hiển nhiên đúng. Do vậy, theo định nghĩa, $A \subseteq A$. Điều phải chứng minh.

\stepcounter{subexercise}
\arabic{subexercise}. Theo định nghĩa tập con,
\begin{equation*}
   \begin{cases}
      A \subseteq B \iff \forall z, (z \in A \implies z \in B) \\
      B \subseteq A \iff \forall z, (z \in B \implies z \in A)
   \end{cases}.
\end{equation*}
Do đó,
\begin{align*}
   A \subseteq B \land B \subseteq A & \iff \forall z, (z \in A \implies z \in B) \land \forall z, (z \in B \implies z \in A);    \\
   A \subseteq B \land B \subseteq A & \iff \forall z, \left( (z \in A \implies z \in B) \land (z \in B \implies z \in A)\right); \\
   A \subseteq B \land B \subseteq A & \iff \forall z, \left(z \in A \iff z \in B\right);                                         \\
   A \subseteq B \land B \subseteq A & \iff A = B.
\end{align*}
Điều phải chứng minh.

\stepcounter{subexercise}
\arabic{subexercise}. Khi biết $\overline{P \iff Q} \iff (P \land \neg Q) \lor (\neg P \land Q)$, biến đổi cơ bản cho chúng ta:
\begin{align*}
   A \neq B & \iff \overline{A = B};                                                                                \\
   A \neq B & \iff \overline{\forall z, z \in A \iff z \in B};                                                      \\
   A \neq B & \iff \exists z, \overline{z \in A \iff z \in B};                                                      \\
   A \neq B & \iff \exists z, \left( z \in A \land z \notin B \right) \lor \left( z \notin A \land z \in B \right).
\end{align*}
Điều phải chứng minh.

\stepcounter{subexercise}
\arabic{subexercise}.
\begin{equation*}
   \begin{cases}
      A \subseteq B \iff \forall z, (z \in A \implies z \in B) \\
      B \subseteq C \iff \forall z, (z \in B \implies z \in C)
   \end{cases}.
\end{equation*}
Do đó,
\begin{align*}
   A \subseteq B \land B \subseteq C & \iff \forall z, (z \in A \implies z \in B) \land \forall z, (z \in B \implies z \in C);    \\
   A \subseteq B \land B \subseteq C & \iff \forall z, \left( (z \in A \implies z \in B) \land (z \in B \implies z \in C)\right); \\
   A \subseteq B \land B \subseteq C & \implies \forall z, \left(z \in A \implies z \in C\right);                                 \\
   A \subseteq B \land B \subseteq C & \implies A \subseteq C.
\end{align*}
Điều phải chứng minh.

\exercise Cho vị từ $V()$. Chứng minh rằng 
\begin{equation*}
   \exists ! x, V(x) \implies \exists x, V(x) \land \forall y, \forall z, (V(y) \land V(z) \implies y = z).
\end{equation*}
($x$, $y$, $z$ là tập hợp.)

\solution

Giả sử $\exists ! x, V(x)$, tương đương với $\exists x, (V(x) \land \forall y, (V(y) \implies x = y))$. Cụ thể hóa $x$ thành $X$, chúng ta có $V(X)$ và $\forall y, (V(y) \implies X = y)$. Đổi hạng thức $y$ thành hạng thức $z$, được mệnh đề tương đương $\forall z, (V(z) \implies X = z)$. Bổ sung lượng từ rỗng,
\begin{equation*}
   \begin{cases}
      \forall y, \forall z, (V(y) \implies X = y) \\ 
      \forall y, \forall z, (V(z) \implies X = z)
   \end{cases}.
\end{equation*}
Do lượng từ phổ quát có tính chất phân phối đối với phép hội nên suy ra 
\begin{equation*}
   \forall y, \forall z, \begin{cases}
      V(y) \implies X = y \\
      V(z) \implies X = z
   \end{cases}.
\end{equation*}

Có $\forall y, \forall z, \begin{cases}
   V(y) \land V(z) \implies V(y) \\
   V(y) \land V(z) \implies V(z)
\end{cases}$ theo quy tắc rút gọn. Kết hợp với tính chất bắc cầu của phép kéo theo để có \begin{equation*}
   \forall y, \forall z, \begin{cases}
      V(y) \land V(z) \implies X = y \\
      V(y) \land V(z) \implies X = z
   \end{cases}.
\end{equation*}
Sử dụng tính chất phân phối một lần nữa, \begin{equation*}
   \forall y, \forall z, (V(y) \land V(z) \implies X = y \land X = z).
\end{equation*}
Do tính bắc cầu của phép đồng nhất, nên cuối cùng cũng suy ra được \begin{equation*}
   \forall y, \forall z, (V(y) \land V(z) \implies y = z).
\end{equation*}
Kết hợp với $\exists x, V(x)$ có được do $V(X)$, khi đó \begin{equation*}
   \exists x, V(x) \land \forall y, \forall z, (V(y) \land V(z) \implies y = z).
\end{equation*}
Vì vậy, chúng ta có điều phải chứng minh.



\section{Tiên đề tập rỗng}

\ 

Tiên đề này phát biểu rằng tồn tại một tập hợp mà không có phần tử nào. Trình bày dưới dạng kí hiệu:
$$
\defMath{\exists x, \forall y, y \notin x}.
$$

Tập này được gọi là \defText{tập rỗng} và được kí hiệu là $\defMath{\varnothing}$ hoặc $\defMath{\emptyset}$. Tập này là duy nhất, do nếu giả sử tồn tại hai tập rỗng $\varnothing \neq \varnothing '$. Chúng ta có một số phân tích với phép $\neq$ như sau: Với $x$ và $y$ là hai tập hợp,
\begin{align*}
   x = y &\iff \forall z, (z \in x \iff z \in y); \\
   x \neq y &\iff \overline{\forall z, (z \in x \iff z \in y)}; \\
   x \neq y &\iff \exists z, \overline{(z \in x \iff z \in y)}.
\end{align*}
Sử dụng bảng giá trị chân lí, chúng ta có $\neg(P \iff Q) \iff (P \land \neg Q) \lor (\neg P \land Q)$.

\begin{table}[H]
	\centering
	\caption{Bảng giá trị chân lí của $\neg ( P \iff Q ) \iff ( P \land \neg Q ) \lor ( \neg P \land Q )$.}
	\begin{tabular}{|c|c|cccccccccccccc|}
		\hline
		$P$ & $Q$ & $\neg$ & $(P$ & $\iff$ & $Q)$ & $\iff$ & $(P$ & $\land$ & $\neg$ & $Q)$ & $\lor$ & $(\neg$ & $P$ & $\land$ & $Q)$ \\
		\headerDivider
		Đ & Đ & S & Đ & Đ & Đ & \emphcolor{Đ} & Đ & S & S & Đ & S & S & Đ & S & Đ \\
		Đ & S & Đ & Đ & S & S & \emphcolor{Đ} & Đ & Đ & Đ & S & Đ & S & Đ & S & S \\
		S & Đ & Đ & S & S & Đ & \emphcolor{Đ} & S & S & S & Đ & Đ & Đ & S & Đ & Đ \\
		S & S & S & S & Đ & S & \emphcolor{Đ} & S & S & Đ & S & S & Đ & S & S & S \\
		\hline
	\end{tabular}
\end{table}
Sử dụng kết quả này để có
\begin{equation*}
   x \neq y \iff \exists z, (z \in x \land z \notin y) \lor (z \notin x \land z \in y).
\end{equation*}
Qua đó, từ $\varnothing \neq \varnothing '$, chúng ta cần phải có
$$
   \exists z, (z \in \varnothing \land z \notin \varnothing ') \lor (z \notin \varnothing \land z \in \varnothing ').
$$
Tuy nhiên, $z \in \varnothing$ luôn sai, và $z \in \varnothing '$ cũng luôn sai theo định nghĩa của tập rỗng. Do đó, cả hai phần của mệnh đề trên đều sai, dẫn đến một mâu thuẫn. Từ đó, chúng ta có chứng minh bằng phản chứng và mọi tập rỗng đều bằng nhau.

Nó có một tính chất mà có cảm giác như một nghịch lí: một tập rỗng là tập con của mọi tập hợp. Thật vậy, theo định nghĩa tập con,
$$
\forall y, (\varnothing \subseteq y \iff \forall z, (z \in \varnothing \implies z \in y)).
$$
Tuy nhiên, $z \notin \varnothing$ đúng với mọi $z$, cho nên vế giả thiết của $z \in \varnothing \implies z \in y$ luôn sai. Do đó, mệnh đề kéo theo này luôn đúng. Từ đó, hiển nhiên dẫn đến $\forall y, (\varnothing \subseteq y)$ đúng.

Liên quan một chút đến vấn đề triết học, người ta có thể đặt câu hỏi: ``Cần tối thiểu bao nhiêu phần tử cơ bản để tạo thành một thế giới?''. Với lí thuyết tập hợp, một tập rỗng là đủ. Trong những phần tiếp theo, bạn đọc sẽ thấy rằng tất cả các cấu trúc toán học đều có thể được xây dựng từ tập hợp. Tuy nhiên (và cũng tương đối hiển nhiên), một thế giới chỉ có một tập rỗng thì sẽ không có gì xảy ra cả, và cũng không có ai để nhận thức về nó. Do đó, mặc dù một tập rỗng là đủ để tạo ra một thế giới toán học, nhưng về mặt thực tế vật lí, chúng ta cần nhiều hơn thế.

Quay trở lại về các vấn đề liên quan đến tập hợp, chúng ta đã có điểm khởi đầu --- $\varnothing$. Những tiên đề tiếp theo sẽ cho phép xây dựng các tập hợp phức tạp hơn từ điểm khởi đầu này.

\exercise Cho $x$ là một tập hợp. Chứng minh rằng $x \subseteq \varnothing \iff x = \varnothing$.

\solution

Có thể thấy, $x = \varnothing \implies x \subseteq \varnothing$ là hiển nhiên do tập rỗng là tập con của mọi tập hợp. 

Trong một trường hợp khác, giả sử $x \subseteq \varnothing$. Điều này có nghĩa là $\forall y, (y \in x \implies y \in \varnothing)$. Tuy nhiên, $y \in \varnothing$ luôn sai do tiên đề tập rỗng. Do đó, $$\forall y, y \notin x.$$ Theo định nghĩa tập rỗng và tính duy nhất của tập này, $x = \varnothing$. Vậy,
$$x \subseteq \varnothing \implies x = \varnothing.$$

Từ hai chiều, chúng ta có được điều phải chứng minh.
\section{Tiên đề cặp}

\ 

Tiên đề này cho phép chúng ta tạo ra một tập hợp mới từ hai tập hợp đã cho. Nó phát biểu rằng với mọi tập hợp $x$ và $y$, tồn tại một tập hợp $z$ sao cho $z$ chỉ chứa $x$ và $y$ làm phần tử của nó. Phát biểu dưới dạng ngôn ngữ lô-gích:
\begin{equation*}
    \defMath{\forall x, \forall y, \exists z, \forall w \left( w \in z \iff (w = x \lor w = y) \right)}.
\end{equation*}

Có thể thấy chỉ tồn tại một $z$ duy nhất theo tiên đề ngoại diên như khi chúng ta đã chứng minh về tính duy nhất của $\varnothing$. Giả sử tồn tại $z_1$ và $z_2$ thỏa mãn tiên đề cặp. Theo chính tiên đề này:
\begin{equation*}
    \begin{cases}
        a \in z_1 \iff (a = x \lor a = y) \\
        a \in z_2 \iff (a = x \lor a = y)
    \end{cases}.
\end{equation*}
Theo tính chất bắc cầu của phép đẳng giá, chúng ta có:
$$
a \in z_1 \iff a \in z_2
$$
đúng với mọi $a$. Theo tiên đề ngoại diên, $z_1 = z_2$.

Kí hiệu tập này là $\defMath{\{x; y\}}$. Với một lập luận tương đương, dễ dàng thấy rằng $\defMath{\{x; y\} = \{y; x\}}$. Khi này, chúng ta có một \defText{đôi không thứ tự} hay \defText{cặp không thứ tự}. Một điểm nữa cần phải để ý, tiên đề không yêu cầu $x \neq y$, nên chúng ta hoàn toàn có thể xây dựng tập $\{x; x\}$ mà chỉ có một phần tử $x$. Đây là \defText{tập hợp đơn}, và có thể viết tắt tập này là $\defMath{\{x\}}$.

Từ một đôi không thứ tự, chúng ta có thể xây dựng được \defText{đôi có thứ tự} (hoặc \defText{cặp có thứ tự}) theo định nghĩa của Ku-ra-tốp-xki\footnote{Kazimierz Kuratowski (1896 -- 1980).} vào năm 1921: $$\forall x, \forall y,\defMath{(x; y) = \{\{x; y\}; \{x\}\}}.$$
Ku-ra-tốp-xki không phải là người đầu tiên đưa định nghĩa về đôi có thứ tự, nhưng là định nghĩa này trở nên phổ biến do tính đơn giản và hiệu quả của nó. Cụ thể, định nghĩa này đảm bảo rằng:
\begin{quotation}
    Với mọi tập hợp $x$, $y$, $u$, $v$, có $(x; y) = (u; v)$ nếu và chỉ nếu $x = u$ và $y = v$.
\end{quotation}
Để chứng minh điều này, chúng ta chia làm ba trường hợp:

\textcolor{colorEmphasisCyan}{Trường hợp 1 --- $x = y$}: Khi này, $$(x; y) = (x; x) = \{\{x; x\}; \{x\}\} = \{\{x\}; \{x\}\} = \{\{x\}\}.$$
Ngoài ra, còn có $(x; y) = (u; v)$ và định nghĩa $(u; v) = \{\{u; v\}; \{u\}\}$. Điều này dẫn đến 
$$
\{\{x\}\} = \{\{u; v\}; \{u\}\}.
$$
Tiên đề ngoại diên cho $\{x\} = \{u; v\}$. Sử dụng tiên đề ngoại diên một lần nữa để có $x = u = v$. Theo tính chất bắc cầu, chúng ta có $x = u = y = v$.

\textcolor{colorEmphasis}{Trường hợp 2 --- $u = v$}: Lập luận tương tự như trường hợp 1 để có điều tương đương.

\textcolor{colorEmphasisGreen}{Trường hợp 3 --- $x \neq y$ và $u \neq v$}: Do $\{x\} \in (x; y)$ và $(x; y) = (u; v)$ nên $\{x\} \in (u; v)$. Từ đó, $\{x\} = \{u; v\}$ hoặc $\{x\} = \{u\}$ theo tiên đề hợp. Tuy nhiên, $\{x\} = \{u; v\}$ dẫn đến $x = u = v$, mâu thuẫn với giả thiết $u \neq v$. Vậy $\{x\} = \{u\}$, tức là $x = u$ theo tiên đề ngoại diên.

$(x; y) = (u; v)$ cũng cho chúng ta $\{x; y\} \in (u; v)$, hay $\{x; y\} = \{u; v\}$ hoặc $\{x; y\} = \{u\}$. Nhưng nếu $\{x; y\} = \{u\}$, thì, sử dụng tiên đề ngoại diên, $x = y = u$, mâu thuẫn với giả thiết $x \neq y$. Vậy $\{x; y\} = \{u; v\}$, tức là $y = u$ hoặc $y = v$; nhưng nếu $y = u$ thì $x = y = u$, mâu thuẫn. Vậy $y = v$. 

Kết hợp ba trường hợp, chúng ta có điều phải chứng minh.

\exercise Chứng minh các mệnh đề sau:
\begin{multicols}{2}
   \begin{enumerate}
      \item $\forall x, \forall y, \left(x = \{y\} \implies x \ni y\right)$;
      \item $\forall x, \forall y, \left(x = y \iff \{x\} = \{y\}\right)$;
      \item $\forall x, \forall y, \left(x \in y \iff \{x\} \subseteq y\right)$;
      \item $\forall x, \forall y, \left(\{x; y\} = \{x\} \iff x = y \right)$.
   \end{enumerate}
\end{multicols}

\solution

\setcounter{subexercise}{1}
\arabic{subexercise}. Giả sử $x = \{y\}$. Theo tính phản xạ của phép bằng nhau, $y = y$ với mọi $y$. Do vậy, $y \in \{y\}$. Khi đó, có $x = \{y\} \land y \in \{y\}$, sử dụng nguyên lí không thể phân biệt của các đối tượng đồng nhất, nên $y \in x$.

Do đó, với mọi $x$ và $y$, $x = \{y\} \implies x \ni y$. Điều phải chứng minh.

\stepcounter{subexercise}
\arabic{subexercise}. Trước hết, coi $x = y$ là đúng. Với mọi $z$, theo định nghĩa và tiên đề hợp thì
\begin{equation*}
   z \in \{x\} \iff z = x.
\end{equation*}
Tương tự, $z \in \{y\} \iff z = y$. Khi chúng ta kèm theo $x = y$, cần phải có $z = x \iff z = y$. Theo tính chất bắc cầu,
\begin{equation*}
   \forall z, z \in \{x\} \iff z \in \{y\}.
\end{equation*}
Vậy, $\forall x, \forall y, \left(x = y \implies \{x\} = \{y\}\right)$.

Ngược lại, coi $\{x\} = \{y\}$ đúng, dẫn đến $\forall w, \left(w \in \{x\} \iff w \in \{y\}\right)$. Ngoài ra, tương tự như trước, chúng ta cũng có khẳng định sau đúng với mọi $w$:
\begin{equation*}
   \begin{cases}
      w \in \{x\} \iff w = x \\
      w \in \{y\} \iff w = y
   \end{cases}.
\end{equation*}
Bắc cầu để được $\forall w, (w = x \iff w = y)$. Cụ thể hóa $w$ bởi $x$, chúng ta được $x = x \iff x = y$. Mà $x = x$ luôn đúng do tính phản xạ nên $x = y$. Vậy
\begin{equation*}
   \forall x, \forall y, \left(\{x\} = \{y\} \implies x = y\right).
\end{equation*}

Kết hợp hai chiều, điều phải chứng minh.

\stepcounter{subexercise}
\arabic{subexercise}. Thứ nhất, xét trường hợp đã có $\{x\} \subseteq y$. Khi này, theo định nghĩa tập con,
\begin{equation*}
   \forall w, \left(w \in \{x\} \implies w \in y\right).\\
\end{equation*}
Đổi $w$ bởi $x$, chúng ta có
\begin{equation*}
   x \in \{x\} \implies x \in y.
\end{equation*}
Mà $x \in \{x\}$ là hiển nhiên, nên $x \in y$.

Thứ hai, xét khi chúng ta có $x \in y$. Xét một phần tử $v$ thuộc $\{x\}$ bất kì. Dễ thấy $v = x$, và do đó $v \in y$. Vậy,
\begin{equation*}
   \forall v, v \in \{x\} \implies v \in y.
\end{equation*}
Dẫn đến, $\{x\} \subseteq y$.

Chúng ta đã chứng minh hai chiều --- $x \in y \impliedby \{x\} \subseteq y$ và $x \in y \implies \{x\} \subseteq y$ --- với mọi $x$ và $y$. Từ định nghĩa phép đẳng giá, chúng ta có điều phải chứng minh.
\section{Tiên đề hợp}

\

Như đã khẳng định, nếu $x$ là một tập hợp thì bên trong $x$ có thể chứa các tập hợp khác. Tiên đề hợp phát biểu rằng tồn tại một tập hợp mà nó chứa tất cả và chỉ chứa các phần tử của các tập hợp trong $x$. Viết dưới dạng kí hiệu:
\begin{equation*}
   \defMath{\forall x, \exists u, \forall y, \left( y \in u \iff \exists z (z \in x \land y \in z) \right)}.
\end{equation*}

\begin{figure}[H]
   \begin{center}
      \begin{tikzpicture}[>=stealth]
         \draw[thick] (-3,0) ellipse (2.5cm and 3cm);
         \node at (-3, 3.4) {Tập hợp $x = \{A, B\}$};

         % Vẽ tập con X bên trong A
         \draw[thick] (-3.8, 0.8) ellipse (0.8cm and 1.2cm);
         \node at (-3.8, 2.3) {$A$};

         % Vẽ tập con Y bên trong A
         \draw[thick] (-2.2, -0.8) ellipse (0.8cm and 1.2cm);
         \node at (-2.2, -2.3) {$B$};

         % Các phần tử trong tập X
         \node[element, label=left:$z_1$] (z1_left) at (-3.8, 1.3) {};
         \node[element, label=left:$z_2$] (z2_left) at (-3.8, 0.3) {};

         % Các phần tử trong tập Y
         \node[element, label=left:$z_3$] (z3_left) at (-2.2, -0.3) {};
         \node[element, label=left:$z_4$] (z4_left) at (-2.2, -1.3) {};

         % --- MŨI TÊN CHUYỂN ĐỔI ---
         \draw[direction arrow, ultra thick, colorEmphasis] (0,0) -- (2,0) node[midway, above] {$\bigcup x$};

         % --- VẾ PHẢI: TẬP HỢP \left(\bigcup A\right) ---
         % Vẽ tập hợp kết quả
         \draw[thick] (4.5,0) ellipse (2cm and 2.8cm);
         \node at (4.5, 3.2) {Tập hợp $\bigcup x$};

         % Các phần tử được "giải phóng" ra tập lớn
         \node[element, label=right:$z_1$] (z1_right) at (4.2, 1.5) {};
         \node[element, label=right:$z_2$] (z2_right) at (4.8, 0.5) {};
         \node[element, label=right:$z_3$] (z3_right) at (4.0, -0.5) {};
         \node[element, label=right:$z_4$] (z4_right) at (4.5, -1.5) {};

         % Nối mũi tên minh họa sự di chuyển của phần tử (giống style của bạn)
         \draw[guide arrow, colorEmphasisCyan, dashed] (z1_left) to[bend left=15] (z1_right);
         \draw[guide arrow, colorEmphasisCyan, dashed] (z3_left) to[bend right=15] (z3_right);

      \end{tikzpicture}
   \end{center}
   \caption{Minh họa tiên đề hợp (các phần tử của $A$ và $B$ trở thành phần tử trực tiếp của $\bigcup x$).}
\end{figure}

Dễ dàng nhận ra rằng tập này là duy nhất. Giả sử có hai tập $u_1$ và $u_2$ đều chứa tất cả và chỉ chứa các phần tử của các tập hợp trong $x$, khi đó ta có:
\begin{equation*}
   \begin{cases}
      \forall y \left( y \in u_1 \iff \exists z (z \in x \land y \in z) \right) \\
      \forall y \left( y \in u_2 \iff \exists z (z \in x \land y \in z) \right)
   \end{cases}.
\end{equation*}
Từ đó suy ra $\forall y, (y \in u_1 \iff y \in u_2)$, tức là $u_1 = u_2$.

Tập hợp này được gọi là \defText{hợp} của $x$ và kí hiệu là $$\defMath{\bigcup x}\text{ hoặc }\defMath{\bigcup_{A \in x} A}.$$

Đặt $x = \{A; B\}$. Khi này, mọi phần tử của A và B đều là phần tử của $\bigcup x$ và ngược lại. Chúng ta có thể hiểu $\bigcup x$ như là hợp của hai tập hợp A và B, tức định nghĩa tương đương \begin{equation*}
   \defMath{A \cup B = \bigcup \{A; B\}}.
\end{equation*}

\exercise Có các tập hợp $A$, $B$ và $C$. Chứng minh rằng tồn tại duy nhất một tập hợp $S$ thỏa mãn $\forall x, (x \in S \iff x = A \lor x = B \lor x = C)$.

\solution

Áp dụng liên tiếp tiên đề cặp, chúng ta có thể xây dựng được tập $\big\{\{A; B\}; \{C\}\big\}$. Thêm tiên đề hợp, chúng ta có $\bigcup \big\{\{A; B\}; \{C\}\big\}$. Chúng ta sẽ chứng minh $$S = \bigcup \big\{\{A; B\}; \{C\}\big\}$$ và đây là tập duy nhất thỏa mãn.

Lấy tùy ý $x \in S$. Theo định nghĩa của hợp, tồn tại tập hợp $y \in \big\{\{A; B\}; \{C\}\big\}$ sao cho $x \in y$. Khi này, $\left[\begin{array}{l}
      y = \{A; B\} \\
      y = \{C\}
   \end{array}\right.$. Xét hai trường hợp:
\begin{itemize}
   \item \textcolor{colorEmphasis}{Trường hợp $y = \{A; B\}$}: Chúng ta có $x \in \{A; B\}$, dẫn đến $x = A$ hoặc $x = B$.
   \item \textcolor{colorEmphasisCyan}{Trường hợp $y = \{C\}$}: Chúng ta có $x \in \{C\}$, dẫn đến $x = C$.
\end{itemize}
Cả hai trường hợp đều cho kết quả $x = A$ hoặc $x = B$ hoặc $x = C$. Do đó, $\forall x, (x \in S \implies x = A \lor x = B \lor x = C)$.

Ngược lại, lấy tùy ý $z$ sao cho $\left[\begin{array}{l}
      z = A \\
      z = B \\
      z = C
   \end{array}\right.$. Xét ba trường hợp:
\begin{itemize}
   \item \textcolor{colorEmphasis}{Trường hợp $z = A$}: Chúng ta có $A \in \{A; B\}$. Ngoài ra, $\{A; B\} \in \big\{\{A; B\}; \{C\}\big\}$, dẫn đến $A \in \bigcup \big\{\{A; B\}; \{C\}\big\}$. Suy ra $z \in S$.
   \item \textcolor{colorEmphasisCyan}{Trường hợp $z = B$}: Tương tự với trường hợp trước, chúng ta cũng có $z \in S$.
   \item \textcolor{colorEmphasisGreen}{Trường hợp $z = C$}: Chúng ta có $C \in \{C\}$ và $\{C\} \in \big\{\{A; B\}; \{C\}\big\}$, dẫn đến $C \in \bigcup \big\{\{A; B\}; \{C\}\big\}$. Suy ra $z \in S$.
\end{itemize}
Tất cả ba trường hợp đều cho kết quả $z \in S$. Do đó, $\forall z, (z = A \lor z = B \lor z = C \implies z \in S)$.

Từ hai điều trên, chúng ta được $\forall x, (x \in S \iff x = A \lor x = B \lor x = C)$.

Tiếp theo, giả sử tồn tại một tập hợp $S'$ khác cũng có tính chất $\forall x, (x \in S' \iff x = A \lor x = B \lor x = C)$. Từ đó suy ra $\forall x, (x \in S \iff x \in S')$, dẫn đến $S = S'$. Điều này chứng tỏ $S$ là tập duy nhất. Chúng ta có điều phải chứng minh.

Chúng ta đặt $S$ bởi một kí hiệu quen thuộc hơn là $\{A; B; C\}$. Ngoài ra, chúng ta có thể mở rộng định nghĩa này cho nhiều phần tử hơn. Cụ thể, giả sử chúng ta có một tập $\{A_1; A_2; \cdots; A_n\}$. Chúng ta có thể thêm một phần tử vào tập này bằng cách đặt $$\defMath{\{A_1; A_2; \cdots; A_n; A_{n+1}\} = \bigcup \left\{\{A_1; A_2; \cdots; A_n\}; \{A_{n+1}\}\right\}}.$$

\exercise Cho $A$, $B$, $C$ là các tập hợp. Chứng minh rằng
\begin{multicols}{2}
   \begin{enumerate}
      \item $\bigcup \{A\} = A$;
      \item $\bigcup \varnothing = \varnothing$;
      \item $A \cup B \supseteq A$;
      \item $A \cup B = B \iff A \subseteq B$;
      \item $A \cup B = B \cup A$;
      \item $\bigcup (A; B) = \{A; B\}$;
      \item $A \subseteq B \implies \bigcup A \subseteq \bigcup B$;
      \item $A \subseteq B \implies A \cup C \subseteq B \cup C$;
      \item $A = B \implies \begin{cases}
                  \bigcup A = \bigcup B \\
                  A \cup C = B \cup C
               \end{cases}$;
      \item $(A \cup B) \cup C = A \cup (B \cup C) = \bigcup \{A; B; C\}$;
      \item $\left(\bigcup A\right) \cup \left(\bigcup B\right) = \bigcup(A \cup B)$.
   \end{enumerate}
\end{multicols}

\solution

\setcounter{subexercise}{1}
\arabic{subexercise}. Đặt $u = \bigcup \{A\}$. Gọi $x$ là một phần tử bất kì thuộc $A$. Hiển nhiên $A \in \{A\}$. Kết hợp với $x \in A$ để có
$$
   \exists z, (z \in \{A\} \land x \in z).
$$
Điều trên là đúng với mọi $x$. Từ đó, suy ra $\forall x, (x \in A \implies x \in u)$, dẫn đến $A \subseteq u$.

Ngược lại, gọi $y$ là một phần tử bất kì thuộc $u$. Theo định nghĩa của tiên đề hợp, chúng ta có
$$
   \exists z, (z \in \{A\} \land y \in z).
$$
Tuy nhiên, $\{A\}$ là một tập hợp chỉ có một phần tử duy nhất là $A$. Do đó, $z = A$ và $y \in A$. Tóm lại, với mọi $y$, nếu $y \in u$ thì $y \in A$, cho chúng ta $u \subseteq A$.

Từ hai điều trên, chúng ta có điều phải chứng minh.

\stepcounter{subexercise}
\arabic{subexercise}. Chúng ta chứng minh bằng phương pháp phản chứng. Giả sử $\bigcup \varnothing \neq \varnothing$, đã có $x \notin \varnothing$ với mọi $x$, cho nên phải tồn tại ít nhất một phần tử $x \in \bigcup \varnothing$. Theo định nghĩa của hợp, sẽ tồn tại một $y \in \varnothing$ sao cho $x \in y$. Tuy nhiên điều này là vô lí, vì tập hợp rỗng $\varnothing$ không có chứa phần tử nào. Sự mâu thuẫn này chứng tỏ giả thiết phản chứng là sai. Vậy $\bigcup \varnothing = \varnothing$. Chúng ta có điều phải chứng minh.

\stepcounter{subexercise}
\arabic{subexercise}. Gọi $x \in A$. Nhắc lại rằng, $A \cup B = \bigcup \{A; B\}$. Khi đó, có $A \in \{A; B\}$ và giả sử ban đầu của $x$ để có $x \in A \cup B$ với mọi $x$. Vậy $A \cup B \supseteq A$. Điều phải chứng minh.

\stepcounter{subexercise}
\arabic{subexercise}. Giả sử $A \cup B = B$. Lấy tùy ý $x \in A$. Áp dụng kết quả của câu trước, chúng ta có $A \subseteq A \cup B$, từ đó suy ra $x \in A \cup B$. Vì giả thiết $A \cup B = B$, nên điều này lập tức kéo theo $x \in B$. Mọi phần tử $x \in A$ đều thuộc $B$, dẫn đến $A \subseteq B$.

Chúng ta giờ chứng minh theo chiều ngược lại. Giả sử $A \subseteq B$.
\begin{itemize}
   \item Lấy tùy ý $a \in A \cup B$. Theo định nghĩa, $a \in A$ hoặc $a \in B$.
         \begin{itemize}
            \item \textcolor{colorEmphasis}{Trường hợp $a \in A$}: Vì $A \subseteq B$ theo giả thiết, $a \in B$.
            \item \textcolor{colorEmphasisCyan}{Trường hợp $a \in B$}: Trực tiếp có được $a \in B$.
         \end{itemize}
         Cả hai trường hợp đều cho kết quả $a \in B$. Do đó, $A \cup B \subseteq B$.
   \item Lấy tùy ý $b \in B$, mà $B \subseteq A \cup B$ dễ dàng có được từ phần trước của bài tập này, suy ra $b \in A \cup B$. Do đó $B \subseteq A \cup B$.
\end{itemize}
Từ hai điều trên, suy ra $A \cup B = B$.

Kết hợp hai chiều, chúng ta nhận được $A \cup B = B \iff A \subseteq B$. Điều phải chứng minh.

\stepcounter{subexercise}
\arabic{subexercise}. Từ định nghĩa, dễ dàng có được điều phải chứng minh do
$$ A \cup B = \bigcup \{A; B\} = \bigcup \{B; A\} = B \cup A.$$

\stepcounter{subexercise}
\arabic{subexercise}. Trước tiên, nhắc lại định nghĩa đôi có thứ tự: $(A; B) = \big\{ \{A\}; \{A; B\} \big\}$. Lấy $x \in \bigcup (A; B)$. Theo định nghĩa của hợp, tồn tại tập hợp $y \in \big\{ \{A\}; \{A; B\} \big\}$ sao cho $x \in y$. Do đó, $y = \{A\}$ hoặc $y = \{A; B\}$.
\begin{itemize}
   \item \textcolor{colorEmphasis}{Trường hợp $y = \{A\}$}: Chúng ta có $x \in \{A\}$ kéo theo $x = A$. Dẫn đến $x \in \{A; B\}$.
   \item \textcolor{colorEmphasisCyan}{Trường hợp $y = \{A; B\}$}: Chúng ta có ngay $x \in \{A; B\}$.
\end{itemize}
Cả hai trường hợp đều cho $\bigcup (A; B) \subseteq \{A; B\}$.

Ngược lại, giả sử $z \in \{A; B\}$. Có $\{A; B\} \in \big\{\{A\}; \{A; B\}\big\}$, nên $z \in \bigcup (A; B)$. Do vậy, $\{A; B\} \subseteq \bigcup (A; B)$.

Từ hai điều trên, chúng ta có điều phải chứng minh.

\stepcounter{subexercise}
\arabic{subexercise}. Giả sử $A \subseteq B$. Lấy tùy ý $x \in \bigcup A$. Theo định nghĩa, tồn tại $y \in A$ sao cho $x \in y$. Vì $A \subseteq B$, nên $y \in B$. Kết hợp với $x \in y$, chúng ta có $\exists z, (z \in B \land x \in z)$, dẫn đến $x \in \bigcup B$. Vậy $\bigcup A \subseteq \bigcup B$. Điều phải chứng minh.

\stepcounter{subexercise}
\arabic{subexercise}. Một lần nữa, giả sử $A \subseteq B$. Lấy tùy ý $x \in A \cup C$, theo định nghĩa, tồn tại tập hợp $y \in \{A; C\}$ sao cho $x \in y$. Khi này, $\left[\begin{array}{l}
      y = A \\
      y = C
   \end{array}\right.$.
\begin{itemize}
   \item \textcolor{colorEmphasis}{Trường hợp $y = A$}: Chúng ta có $x \in A$. Vì $A \subseteq B$, suy ra $x \in B$. Lại có $B \in \{B; C\}$, nên kéo theo $x$ là một phần tử của $B \cup C$.
   \item \textcolor{colorEmphasisCyan}{Trường hợp $y = C$}: Chúng ta có $x \in C$. Tương tự, có $x \in B \cup C$.
\end{itemize}
Cả hai trường hợp đều cho kết quả $x \in B \cup C$. Do đó, $A \cup C \subseteq B \cup C$.

Vậy $A \subseteq B \implies A \cup C \subseteq B \cup C$. Điều phải chứng minh.

\stepcounter{subexercise}
\arabic{subexercise}. Phần bài tập chỉ mang tính chất bổ sung cho đầy đủ kiến thức. Nếu như bạn đọc đã chứng minh được hai phần trước thì bạn đọc sẽ cảm nhận rằng phần này là một điều hiển nhiên. Có $A = B \iff A \subseteq B \land B \subseteq A$. Từ hai phần trước, chúng ta có
\begin{equation*}
   A \subseteq B \implies \begin{cases}
      \bigcup A \subseteq \bigcup B \\
      A \cup C \subseteq B \cup C
   \end{cases}.
\end{equation*}
Tương tự, $B \subseteq A \implies \begin{cases}
      \bigcup B \subseteq \bigcup A \\
      B \cup C \subseteq A \cup C
   \end{cases}$.
Kết hợp hai điều trên, chúng ta có
\begin{align*}
   A \subseteq B \land B \subseteq A & \implies \begin{cases}
                                                   \bigcup A \subseteq \left(\bigcup B\right) \land \left(\bigcup B\right) \subseteq \bigcup A \\
                                                   A \cup C \subseteq B \cup C \land B \cup C \subseteq A \cup C
                                                \end{cases} \\
   A = B                             & \implies \begin{cases}
                                                   \bigcup A = \left(\bigcup B\right) \\
                                                   A \cup C = B \cup C
                                                \end{cases}.
\end{align*}
Điều phải chứng minh.

\stepcounter{subexercise}
\arabic{subexercise}. Trước tiên, chúng ta đi chứng minh $(A \cup B) \cup C = \bigcup \{A; B; C\}$.

Giả sử $x \in (A \cup B) \cup C$. Theo định nghĩa, $x \in A \cup B$ hoặc $x \in C$. Nếu $x \in A \cup B$, thì tiếp tục áp dụng định nghĩa để có $x \in A$ hoặc $x \in B$. Tóm lại, $x \in A$, $x \in B$ hoặc $x \in C$; có nghĩa là tồn tại tập hợp $y \in \{A; B; C\}$ sao cho $x \in y$. Điều này chứng tỏ $x \in \bigcup \{A; B; C\}$, và kéo theo $(A \cup B) \cup C \subseteq \bigcup \{A; B; C\}$.

Ngược lại, lấy tùy ý $z \in \bigcup \{A; B; C\}$. Theo định nghĩa, tồn tại tập hợp $w \in \{A; B; C\}$ sao cho $z \in w$. Khi đó, $\left[\begin{array}{l}
      w = A \\
      w = B \\
   w = C
   \end{array}\right.$. Xét ba trường hợp:
\begin{itemize}
   \item \textcolor{colorEmphasis}{Trường hợp $w = A$}: Chúng ta có $z \in A$. Từ đó suy ra $z \in A \cup B$, dẫn đến $z \in (A \cup B) \cup C$.
   \item \textcolor{colorEmphasisCyan}{Trường hợp $w = B$}: Chúng ta có $z \in B$. Tương tự, cũng có $z \in A \cup B$, kéo theo $z \in (A \cup B) \cup C$.
   \item \textcolor{colorEmphasisGreen}{Trường hợp $w = C$}: Chúng ta có $z \in C$. Từ đó trực tiếp suy ra $z \in (A \cup B) \cup C$.
\end{itemize}
Tất cả ba trường hợp đều cho kết quả $z \in (A \cup B) \cup C$. Do đó, $\bigcup \{A; B; C\} \subseteq (A \cup B) \cup C$.

Từ hai điều trên, suy ra $(A \cup B) \cup C = \bigcup \{A; B; C\}$.

Áp dụng lập luận biến đổi hoàn toàn tương tự, chúng ta cũng có được $A \cup (B \cup C) = \bigcup \{A; B; C\}$. Do đó, kéo theo $(A \cup B) \cup C = A \cup (B \cup C) = \bigcup \{A; B; C\}$. Chúng ta có điều phải chứng minh.

\stepcounter{subexercise}
\arabic{subexercise}. Với mọi $x$, 
\begin{align*}
   x \in \left(\bigcup A\right) \cup \left(\bigcup B\right) &\iff x \in \bigcup A \lor x \in \bigcup B; \\
   x \in \left(\bigcup A\right) \cup \left(\bigcup B\right) &\iff \left[\begin{array}{l}
      \exists z, (z \in A \land x \in z) \\
      \exists z, (z \in B \land x \in z)
   \end{array}\right. \equationexplanation{tiên đề hợp}; \\
   x \in \left(\bigcup A\right) \cup \left(\bigcup B\right) &\iff \exists z, (z \in A \land x \in z \lor z \in B \land x \in z); \\
   x \in \left(\bigcup A\right) \cup \left(\bigcup B\right) &\iff \exists z, \left((z \in A \lor z \in B) \land x \in z\right); \\
   x \in \left(\bigcup A\right) \cup \left(\bigcup B\right) &\iff \exists z, (z \in A \cup B \land x \in z); \\
   x \in \left(\bigcup A\right) \cup \left(\bigcup B\right) &\iff x \in \bigcup (A \cup B).
\end{align*}
Vậy $\left(\bigcup A\right) \cup \left(\bigcup B\right) = \bigcup (A \cup B)$. Điều phải chứng minh.

\exercise Cho tập hợp $x$ thỏa mãn $A \cup x = A$ với mọi tập hợp $A$. Chứng minh rằng $x = \varnothing$.

\solution

Cụ thể hóa mệnh đề $\forall A, A \cup x = A$ bằng việc thay $A$ bởi $\varnothing$, chúng ta có
\begin{equation*}
   \varnothing \cup x = \varnothing.
\end{equation*}
Do $\varnothing \subseteq x$ nên $\varnothing \cup x = x$. Vì vậy, 
\begin{equation*}
   x = \varnothing.
\end{equation*}
Điều phải chứng minh.
\section{Tiên đề tập lũy thừa}

\

Tiên đề tập lũy thừa phát biểu rằng với mọi tập hợp $x$, tồn tại một tập hợp $p$ mà nó chứa tất cả và chỉ chứa các phần tử là các tập hợp con của $x$. Viết dưới dạng kí hiệu:
\begin{equation*}
    \defMath{\forall x, \exists p, \forall y, \left( y \in p \iff y \subseteq x \right)}.
\end{equation*}
Tập $p$ thỏa mãn điều này được gọi là \defText{tập lũy thừa} của $x$, kí hiệu là $\defMath{\mathfrak{P}(x)}$, $\defMath{\mathcal{P}(x)}$, hay $\defMath{2^x}$\footnote{Cần phải để ý, các kí hiệu ở đây là một sự lạm dụng. Kí hiệu $\mathfrak{P}(x)$ không phải để biểu diễn vị từ, và $2^x$ không biểu diễn phép lũy thừa.}.

Cho trước tập hợp $x$, giả sử tồn tại hai tập lũy thừa của $x$ là $\mathfrak{P}(x)$ và $\mathfrak{B}(x)$. Chúng ta sẽ chứng minh rằng $\mathfrak{P}(x) = \mathfrak{B}(x)$. Giống như các chứng minh trước, chúng ta xuất phát từ định nghĩa để đưa ra:
\begin{equation*}
    \begin{cases}
        \forall y, \left( y \in \mathfrak{P}(x) \iff y \subseteq x \right) \\
        \forall y, \left( y \in \mathfrak{B}(x) \iff y \subseteq x \right)
    \end{cases}.
\end{equation*}
Từ đó, có được ngay, $\forall y, (y \in \mathfrak{P}(x) \iff y \in \mathfrak{B}(x))$, suy ra $\mathfrak{P}(x) = \mathfrak{B}(x)$.

\exercise Cho $X$ và $Y$ là hai tập hợp. Chứng minh rằng,
\begin{multicols}{2}
    \begin{enumerate}
        \item $\varnothing \in \mathfrak{P}(X)$;
        \item $\mathfrak{P}(\{X; Y\}) = \{\varnothing; \{X\}; \{Y\}; \{X; Y\}\}$;
        \item $X \subseteq Y \implies \mathfrak{P}(X) \subseteq \mathfrak{P}(Y)$;
        \item $\mathfrak{P}(X \cup Y) \supseteq \mathfrak{P}(X) \cup \mathfrak{P}(Y)$.
    \end{enumerate}
\end{multicols}

\solution

\setcounter{subexercise}{1}
\arabic{subexercise}. Có tập rỗng là tập con của mọi tập hợp. Cho nên $\varnothing \subseteq X$. Nhắc lại định nghĩa của tập lũy thừa, $\mathfrak{P}(X)$ là tập chứa tất cả các tập con của $X$. Do đó, $\varnothing \in \mathfrak{P}(X)$. Điều phải chứng minh.

Một hệ quả tự nhiên của tính chất này đó là tập lũy thừa thì không rỗng. Biểu diễn dưới dạng kí hiệu:
$$
    \forall w, \mathfrak{P}(w) \neq \varnothing.
$$

\stepcounter{subexercise}
\arabic{subexercise}. Tập rỗng là tập con của mọi tập hợp, nên $\varnothing \subseteq \{X; Y\}$, dẫn đến $\varnothing \in \mathfrak{P}(\{X; Y\})$.

Theo cách xây dựng của đôi không thứ tự, $X \in \{X; Y\}$. Tập hợp đơn $\{X\}$ chỉ có phần tử là $X$. Do vậy, $\{X\} \subseteq \{X; Y\}$, dẫn đến $\{X\} \in \mathfrak{P}(\{X; Y\})$. Tương tự, $\{Y\} \in \mathfrak{P}(\{X; Y\})$.

Hiển nhiên $\{X; Y\} = \{X; Y\}$. Do vậy, $\{X; Y\} \in \mathfrak{P}(\{X; Y\})$.

Từ ba kết luận trên, chúng ta có $\mathfrak{P}(\{X; Y\}) \supseteq \{\varnothing; \{X\}; \{Y\}; \{X; Y\}\}$.

Mặt khác, với mọi $y \in \mathfrak{P}(\{X; Y\})$, theo định nghĩa có $y \subseteq \{X; Y\}$, hiểu theo một cách khác, $y$ chỉ nhận $X$ và $Y$ làm phần tử. Do mỗi phần tử $X$ và $Y$ chỉ có thể thuộc hoặc không thuộc $y$, nên các khả năng có thể xảy ra bao gồm\footnote{Tác giả sử dụng mở rộng của quy tắc chia trường hợp từ giải tích mệnh đề: ``Cho $P$ và $Q$ là hai mệnh đề, khi này, chúng ta có $P \land Q \lor P \land \neg Q \lor \neg P \land Q \lor \neg P \land \neg Q$.''.}:

\begin{itemize}
    \item \textcolor{colorEmphasis}{$X \in y$ và $Y \in y$}. Điều này cho chúng ta $\{X; Y\} \subseteq y$. Suy ra, $y = \{X; Y\}$.
    \item \textcolor{colorEmphasisCyan}{$X \in y$ và $Y \notin y$}. Gọi $w$ là một phần tử thuộc $y$. Có $\left[\begin{array}{l} w = X \\ w = Y \end{array}\right.$. Nhưng do $Y \notin y$ nên $w \neq Y$. Vậy $w = X$. Tóm lại, chúng ta có
          $$\forall w, \left(w \in y \iff w = X\right)$$
          và điều này cho $y = \{X\}$.
    \item \textcolor{colorEmphasisGreen}{$X \notin y$ và $Y \in y$}. Tương tự, có thể kết luận $y = \{Y\}$.
    \item \textcolor{colorEmphasisPurple}{$X \notin y$ và $Y \notin y$}. Giả sử tồn tại $w \in y$. Khi đó, cần phải có $w = X \lor w = Y$. Nhưng do $X \notin y$ và $Y \notin y$, nên điều này là vô lí. Vậy, $\forall w, w \notin y$. Suy ra, $y = \varnothing$ theo định nghĩa tập rỗng.
\end{itemize}

Qua bốn trường hợp, có $y \in \{\varnothing; \{X\}; \{Y\}; \{X; Y\}\}$. Do đó, $\mathfrak{P}(\{X; Y\}) \subseteq \{\varnothing; \{X\}; \{Y\}; \{X; Y\}\}$.

Kết hợp hai kết quả bao hàm, suy ra $\mathfrak{P}(\{X; Y\}) = \{\varnothing; \{X\}; \{Y\}; \{X; Y\}\}$. Điều phải chứng minh.

\begin{figure}[H]
    \begin{center}
        \begin{tikzpicture}[>=stealth]
            \draw[thick] (-3,0) ellipse (2.1cm and 3cm);
            \node at (-3, 3.3) {Tập hợp $\{a; b\}$};

            \node[circle, fill, inner sep=1.2pt, label=left:$a$] (a_left) at (-4, 0) {};
            \node[circle, fill, inner sep=1.2pt, label=left:$b$] (b_left) at (-2.5, -1.2) {};

            \draw[thick, dashed, colorEmphasis] (-2.5, 1.5) circle (0.4cm);
            \node[colorEmphasis] at (-2.5, 1.5) {$\varnothing$};

            \draw[thick, dashed, colorEmphasisCyan] (-4, 0) ellipse (0.5cm and 0.4cm);
            \node[colorEmphasisCyan] at (-4, 0.7) {$\{a\}$};

            \draw[thick, dashed, colorEmphasisGreen] (-2.5, -1.2) ellipse (0.4cm and 0.5cm);
            \node[colorEmphasisGreen] at (-1.8, -1.2) {$\{b\}$};

            \draw[thick, dashed, colorEmphasisPurple, rotate around={50:(-3.15,-0.45)}] (-3.15, -0.45) ellipse (1cm and 2cm);
            \node[colorEmphasisPurple] at (-4.2, -1.45) {$\{a; b\}$};

            \draw[->, ultra thick] (-0.5, 0) -- (1.9, 0) node[midway, above] {$\mathfrak{P}\left(\{a; b\}\right)$};

            \draw[thick] (5.5,0) ellipse (3.2cm and 3.5cm);
            \node at (5.5, 3.9) {Tập lũy thừa $\mathfrak{P}\left(\{a; b\}\right)$};

            % Tập rỗng
            \draw[thick, colorEmphasis] (5.5, 2) circle (0.7cm);
            \node[colorEmphasis] at (5.5, 2) {$\varnothing$};
            \draw[->, thick, colorEmphasis, dashed, shorten >= 5pt, shorten <= 5pt] (-2.2, 1.7) to[bend left=15] (4.85, 2);

            % Tập {a}
            \draw[thick, colorEmphasisCyan] (3.6, 0) ellipse (0.8cm and 0.6cm);
            \node[circle, fill, inner sep=1.2pt, label=below:$a$] (a_right1) at (3.6, 0) {};
            \node[colorEmphasisCyan] at (3.6, 0.9) {$\{a\}$};
            \draw[->, thick, colorEmphasisCyan, dashed, shorten >= 5pt, shorten <= 5pt] (-3.6, 0) to[bend right=25] (3.2, -0.5);

            % Tập {b}
            \draw[thick, colorEmphasisGreen] (7.4, 0) ellipse (0.6cm and 0.8cm);
            \node[circle, fill, inner sep=1.2pt, label=below:$b$] (b_right1) at (7.4, 0) {};
            \node[colorEmphasisGreen] at (8.3, 0) {$\{b\}$};
            \draw[->, thick, colorEmphasisGreen, dashed, shorten >= 8pt, shorten <= 8pt] (-2.4, -1.5) to[out=-60, in=180] (7, -0.1);

            % Tập {a, b}
            \draw[thick, colorEmphasisPurple, rotate around={-40:(5.5, -2)}] (5.5, -2) ellipse (1.2cm and 0.8cm);
            \node[circle, fill, inner sep=1.2pt, label=left:$a$] (a_right2) at (5.1, -1.7) {};
            \node[circle, fill, inner sep=1.2pt, label=right:$b$] (b_right2) at (5.9, -2.3) {};
            \node[colorEmphasisPurple] at (6.5, -1.2) {$\{a; b\}$};
            \draw[->, thick, colorEmphasisPurple, dashed, shorten >= 8pt, shorten <= 8pt] (-1.7, -1) to[out=50, in=-150] (5.1, -2.6);

        \end{tikzpicture}
    \end{center}
    \caption{Minh họa tiên đề tập lũy thừa với $\mathfrak{P}\left(\{a; b\}\right) = \big\{\varnothing; \{a\}; \{b\}; \{a; b\}\big\}$.}
\end{figure}


\stepcounter{subexercise}
\arabic{subexercise}. Giả sử $X \subseteq Y$. Gọi $w \in \mathfrak{P}(X)$. Theo định nghĩa, $w \subseteq X$. Theo tính chất bắc cầu của quan hệ bao hàm, có $w \subseteq Y$. Do đó, $w \in \mathfrak{P}(Y)$.

Từ đó, có $X \subseteq Y \implies \mathfrak{P}(X) \subseteq \mathfrak{P}(Y)$. Điều phải chứng minh.

\stepcounter{subexercise}
\arabic{subexercise}. Gọi $w$ là một phần tử bất kì thuộc $\mathfrak{P}(X) \cup \mathfrak{P}(Y)$. Theo định nghĩa của phép hợp, có $w \in \mathfrak{P}(X)$ hoặc $w \in \mathfrak{P}(Y)$.
Nếu $w \in \mathfrak{P}(X)$, theo định nghĩa tập lũy thừa, $w \subseteq X$. Mặt khác, $X \subseteq X \cup Y$, nên theo tính chất bắc cầu, $w \subseteq X \cup Y$, dẫn đến $w \in \mathfrak{P}(X \cup Y)$. Chứng minh hoàn toàn tương tự cho trường hợp $w \in \mathfrak{P}(Y)$, chúng ta cũng thu được $w \in \mathfrak{P}(X \cup Y)$.

Vậy, nhìn chung, chúng ta có
$$\forall w, \left( w \in \mathfrak{P}(X) \cup \mathfrak{P}(Y) \implies w \in \mathfrak{P}(X \cup Y) \right)$$
và do đó $\mathfrak{P}(X \cup Y) \supseteq \mathfrak{P}(X) \cup \mathfrak{P}(Y)$. Điều phải chứng minh.


\section{Tiên đề tách}

\

Giả sử $x$ là một tập hợp và $\cauF()$ có nghĩa với toàn bộ các phần tử của $x$. Khi đó, tiên đề tách khẳng định rằng tồn tại một tập hợp $y$ bao gồm tất cả các phần tử $z \in x$ thỏa mãn $\cauF(z)$ và chỉ những phần tử đó. Nói cách khác,
$$
    \defMath{\forall x, \exists y, \forall z, \left( z \in y \iff z \in x \land \textbf{\cauF}(z) \right)}.
$$

\begin{figure}[H]
    \begin{center}
    \begin{tikzpicture}[>=stealth]

    \draw[thick] (-3,0) ellipse (2.2cm and 2.8cm);
    \node at (-3, 3.2) {Tập hợp $x$};

    \node[circle, fill=colorEmphasisCyan, inner sep=\pointSize, label=left:$z_1$] (z1_left) at (-3.7, 0.8) {};
    \node[circle, fill=colorEmphasisCyan, inner sep=\pointSize, label=left:$z_3$] (z3_left) at (-2.5, -0.5) {};

    \node[circle, inner sep=\pointSize, label=left:{$z_2$}] (z2_left) at (-2, 0.8) {};
    \node[circle, inner sep=\pointSize, label=left:{$z_4$}] (z4_left) at (-3.8, -1.2) {};

    \draw[dashed, colorEmphasisCyan, thin] plot [smooth cycle, tension=1] coordinates {
        (-4.6, 0.5)
        (-3.5, 1.8)
        (-3, 0.4)
        (-2, -0.2)
        (-2.5, -1.2) 
    };
    \node at (-3.1, 2.1) {$\textcolor{colorEmphasisCyan}{\cauF(}z\textcolor{colorEmphasisCyan}{)}$};


    \draw[->, ultra thick, colorEmphasis] (0,0) -- (2,0) 
        node[midway, above, align=center] {Tách bởi\\ $\textcolor{colorEmphasisCyan}{\cauF(}z\textcolor{colorEmphasisCyan}{)}$};


    \draw[thick] (4.5,0) ellipse (1.5cm and 2.2cm);
    \node at (4.5, 2.6) {$\{z \in x \mid \textcolor{colorEmphasisCyan}{\cauF(}z\textcolor{colorEmphasisCyan}{)}\}$};

    % Chỉ các phần tử thỏa mãn mới xuất hiện ở đây
    \node[circle, fill=colorEmphasisCyan, inner sep=1.5pt, label=right:$z_1$] (z1_right) at (4.2, 0.8) {};
    \node[circle, fill=colorEmphasisCyan, inner sep=1.5pt, label=right:$z_3$] (z3_right) at (4.8, -0.5) {};


    \draw[->, thick, colorEmphasisCyan, shorten >= 5pt, shorten <= 5pt, dashed] 
        (z1_left) to[bend left=20] (z1_right);

    \draw[->, thick, colorEmphasisCyan, shorten >= 5pt, shorten <= 5pt, dashed] 
        (z3_left) to[bend right=15] (z3_right);

    \end{tikzpicture}
    \end{center}
    \caption{Minh họa tiên đề tách (các phần tử của $x$ thỏa mãn tính chất $\textcolor{colorEmphasisCyan}{\cauF(}z\textcolor{colorEmphasisCyan}{)}$ được tách ra).}
\end{figure}

Tiên đề ngoại diên khẳng định tập hợp $y$ là duy nhất. Thực vậy, giả sử tồn tại hai tập hợp $y_1$ và $y_2$ thỏa mãn điều kiện của tiên đề tách với cùng tập hợp $x$ và vị từ $\cauF()$. Xét một $z$ bất kì. Khi đó, định nghĩa theo tiên đề cho chúng ta
\begin{equation*}
    \begin{cases}
        z \in y_1 \iff z \in x \land \cauF(z) \\
        z \in y_2 \iff z \in x \land \cauF(z)
    \end{cases}.
\end{equation*}
Theo tính bắc cầu của phép đẳng giá, chúng ta thu được
\begin{equation*}
    \forall z, \left( z \in y_1 \iff z \in y_2 \right).
\end{equation*}
Theo tiên đề ngoại diên, $y_1 = y_2$. Gọi tập hợp đó là $y$. Một hệ quả khá hiển nhiên chính là $y \subseteq x$, bởi vì rõ ràng rằng, nếu $z \in y$, theo tiên đề tách, chúng ta phải có $z \in x$.

Trước khi đi vào những vấn đề khác, chúng ta sẽ trình bày về một cách viết tắt các mệnh đề trong giải tích vị từ bậc nhất. Bạn đọc có thể để ý, trong định nghĩa, chúng ta có một mệnh đề dạng
$$
   \forall z, \left(z \in x \implies V(z)\right).
$$
Dưới dạng ngôn ngữ, chúng ta sẽ nói:
\begin{center}
   ``Với mọi, nếu $z$ thuộc $x$ thì $V(z)$.'',
\end{center}
nhưng một cách phát biểu tự nhiên hơn sẽ là:
\begin{center}
   ``Với mọi $z$ thuộc $x$, $V(z)$.''
\end{center}
Từ cách nói đó, chúng ta có cách viết rút gọn với mọi tập hợp $x$ và vị từ $V()$ như sau:
\begin{equation*}
   \defMath{\forall z \in x, V(z)} \iff \forall z, \left(z \in x \implies V(z)\right).
\end{equation*}
Tương tự, với lượng từ tồn tại, chúng ta cũng có cách viết tắt:
\begin{equation*}
   \defMath{\exists z \in x, V(z)} \iff \exists z, \left(z \in x \land V(z)\right).
\end{equation*}

Quay trở lại với tập hợp $y$ trong tiên đề tách, chúng ta không thể kí hiệu bằng việc thêm một kí tự duy nhất như trước được nữa, mà do mỗi tập $y$ đều cần được xác định từ một vị từ $\cauF$, vì thế, chúng ta sẽ kí hiệu tập hợp $y$ như sau:
\begin{equation*}
   y = \defMath{\{ z \in x \mid \cauF(z) \}} = \defMath{\{ z \in x : \cauF(z) \}}.
\end{equation*}
và hệ quả sẽ được kí hiệu là
\begin{equation}
   \{ z \in x \mid \cauF(z) \} \subseteq x.
   \label{eq:li_thuyet_tap_hop:tien_de:tien_de_tach:he_qua_tap_con}
\end{equation}

Liệu bạn đọc còn nhớ về nghịch lí Rất-xen, liên quan đến một tập bao gồm tất cả các tập hợp không chứa chính nó. Có một cách đơn giản để đi vòng qua vấn đề này. Đó chính là không cho nó tồn tại. Để làm như vậy, chúng ta sẽ chứng minh không tồn tại tập hợp $\Omega$ sao cho $\forall x, x \in \Omega$. Thực vậy, giả sử tồn tại tập hợp $\Omega$ như vậy. Khi đó, theo tiên đề tách, tồn tại tập hợp $y = \{x \in \Omega \mid x \notin x \}$. Đặt câu hỏi, $y$ có thuộc $y$ không?
\begin{itemize}
    \item Nếu $y \in y$, thì theo định nghĩa của $y$, chúng ta có $y \notin y$, do $y$ chỉ chứa các tập không chứa chính nó, một mâu thuẫn.
    \item Nếu $y \notin y$, thì tương tự, theo định nghĩa của $y$, chúng ta có $y \in y$, một mâu thuẫn.
\end{itemize}
Nhìn chung, cả hai trường hợp đều đưa ra một điều vô lí. Sử dụng chứng minh bằng phản chứng, chúng ta có thể kết luận rằng không tồn tại tập hợp $\Omega$ sao cho $x \in \Omega$ với mọi $x$.

Giờ là lúc để làm một số ví dụ. Giả sử $x$ và $y$ là hai tập hợp. Nhận thấy rằng, ``$(z) \in y$'' là một vị từ theo hạng thức $z$. Áp dụng tiên đề tách, chúng ta có tập
\begin{equation*}
    \defMath{\{ z \in x \mid z \in y \}}.
\end{equation*}
Đây được gọi là \defText{giao} của $x$ và $y$, kí hiệu $\defMath{x \cap y}$. Nếu $x \cap y = \varnothing$, thì $x$ và $y$ được gọi là \defText{rời nhau}. Ngoài ra, có $z \in y$, thì chúng ta cũng có thể có $z \notin y$, điều đó tự nhiên dẫn đến \defText{hiệu} giữa $x$ và $y$ như sau:
\begin{equation*}
    \defMath{x \setminus y = x - y = \{ z \in x \mid z \notin y \}}.
\end{equation*}
$x \setminus y$ còn được gọi là \defText{phần bù tương đối} của $y$ trong $x$. Thường xuyên, chúng ta sẽ chỉ quan tâm các tính chất giới hạn trong một tập cụ thể, hiểu theo nghĩa, chúng ta sẽ chỉ quan tâm đến tập con của một tập hợp $x$ nào đó. Với giả thiết rằng $y \subseteq x$, chúng ta định nghĩa và kí hiệu \defText{phần bù} của $y$ là $$\defMath{\complement_x(y) = \complement(y) = y^c = \overline{y} = x \setminus y}.$$

Định nghĩa giao có thể được mở rộng lên nhóm nhiều tập hợp. Tiên đề hợp cho phép xây dựng tập $\bigcup x$, và ``$\forall a \in (x), z \in a$'' là một vị từ theo hạng thức $x$. Sử dụng tiên đề tách, chúng ta có tập
{
    \boldmath 
    $$\bigcap x = \bigcap_{a\in x} a = \left\{z \in \bigcup x \;\middle|\; \forall a \in x, z \in a \right\}$$
}%
và gọi nó là giao của các phần tử trong $x$. Bạn đọc cần phải chú ý một điều, để không xảy ra mâu thuẫn, định nghĩa mới phải tương thích với định nghĩa cũ. May mắn thay, $\bigcap \{x; y\} = x \cap y$.

Một phép toán nữa, khá ít dùng, được gọi là \defText{hiệu đối xứng}. Với hai tập hợp $x$ và $y$, hiệu đối xứng của chúng được định nghĩa là
\begin{equation*}
    \defMath{x \vartriangle y = (x \setminus y) \cup (y \setminus x)}.
\end{equation*}

Để kết thúc phần này, tác giả sẽ giới thiệu về cách biểu diễn tập hợp bằng sơ đồ Ven\footnote{John Venn (1834 -- 1923).}. Với hai tập hợp $x$ và $y$, chúng ta có thể biểu diễn chúng bằng hai hình chồng lên nhau, trong đó kết quả của các phép tính trên chúng được tô màu, với một sự minh hoạ cụ thể ở hình \ref{fig:li_thuyet_tap_hop:tien_de:tien_de_tach:ven_phep_tinh}.

\begin{figure}[H]
    \centering
    \begin{tikzpicture}
        \begin{scope}[shift={(0,0)}]
            % Tô màu cả 2 hình tròn
            \fill[color = colorEmphasis] (0,0) circle (0.7);
            \fill[color = colorEmphasis] (0.8,0) circle (0.7);
            
            % % Vẽ đường viền
            \draw (0,0) circle (0.7);
            \draw (0.8,0) circle (0.7);
            
            % Nhãn tên tập hợp (màu trắng)
            \node[above left] at (-0.4, 0.4) {$x$};
            \node[above right] at (1.2, 0.4) {$y$};
            
            % Nhãn phép toán
            \node[below] at (0.4, -0.9) {$x \cup y$};
        \end{scope}

        \begin{scope}[shift={(4,0)}]
            \begin{scope}
                \clip (0,0) circle (0.7);
                \fill[color = colorEmphasis] (0.8,0) circle (0.7);
            \end{scope}
            
            % Vẽ đường viền
            \draw (0,0) circle (0.7);
            \draw (0.8,0) circle (0.7);
            
            % Nhãn
            \node[above left] at (-0.4, 0.4) {$x$};
            \node[above right] at (1.2, 0.4) {$y$};
            \node[below] at (0.4, -0.9) {$x \cap y$};
        \end{scope}

        \begin{scope}[shift={(8, 0)}]
            % Tô màu tập x nhưng bỏ đi phần giao với y
            \begin{scope}
                \clip (0,0) circle (0.7);
                \fill[white] (0.8,0) circle (0.7); % Xóa phần giao bằng nền
            \end{scope}
            % Phủ lại pattern lên phần còn lại của x
            \begin{scope}
                \clip (0,0) circle (0.7);
                \draw (0.8,0) circle (0.7); 
                % Thủ thuật đảo ngược để tô phần x \ y
                \fill[color = colorEmphasis, even odd rule] (0,0) circle (0.7) (0.8,0) circle (0.7);
            \end{scope}
            
            % Vẽ đường viền
            \draw (0,0) circle (0.7);
            \draw (0.8,0) circle (0.7);
            
            % Nhãn
            \node[above left] at (-0.4, 0.4) {$x$};
            \node[above right] at (1.2, 0.4) {$y$};
            \node[below] at (0.4, -0.9) {$x \setminus y$};
        \end{scope}

        \begin{scope}[shift={(12,0)}]
            \fill[color = colorEmphasis, even odd rule] (0,0) circle (0.7) (0.8,0) circle (0.7);
            
            % Vẽ đường viền
            \draw (0,0) circle (0.7);
            \draw (0.8,0) circle (0.7);
            
            % Nhãn
            \node[above left] at (-0.4, 0.4) {$x$};
            \node[above right] at (1.2, 0.4) {$y$};
            \node[below] at (0.4, -0.9) {$x \vartriangle y$};
        \end{scope}

    \end{tikzpicture}
    \caption{Biểu diễn các phép toán tập hợp bằng sơ đồ Ven.}
    \label{fig:li_thuyet_tap_hop:tien_de:tien_de_tach:ven_phep_tinh}
\end{figure}

\exercise Chứng minh rằng, với mọi tập hợp $x$ và $y$, $\bigcap \{x; y\} = x \cap y$.

\solution 

Gọi $z$ là một phần tử trong $\bigcap \{x; y\}$. Khi đó, theo định nghĩa $$\forall a \in \{x; y\}, z \in a.$$
Có $x \in \{x; y\}$ nên $z \in x$, $y \in \{x; y\}$ nên $z \in y$. Do đó, $z \in x \land z \in y$, dẫn đến $z \in x \cup y$. Vậy, $\bigcap \{x; y\} \subseteq x \cap y$.

Ngược lại, giả sử $w \in x \cap y$. Khi này, $w \in x$ và $w \in y$. Nhắc lại rằng $\forall b, b \in \{x; y\} \iff \left[\begin{array}{l}
   b = x \\
   b = y
\end{array}\right.$. Trong cả hai trường hợp đó, $w \in b$. Do vậy, $\forall b \in \{x; y\}, w \in b$. Theo định nghĩa $w \in \bigcap \{x; y\}$. Vậy, $x \cap y \subseteq \bigcap \{x; y\}$.

Kết hợp hai kết quả, chúng ta có $\bigcap \{x; y\} = x \cap y$, theo cách nói khác, điều phải chứng minh.

\exercise Cho $A$, $B$, $C$ là các tập hợp. Chứng minh rằng
\begin{multicols}{2}
   \begin{enumerate}
      \item $\bigcap \{A\} = A$;
      \item $\bigcap \varnothing = \varnothing$
      \item $A \cap B \subseteq A$;
      \item $A \cap B = A \iff A \subseteq B$;
      \item $A \cap B = B \cap A$;
      \item $\bigcap (A; B) = \{A\}$;
      \item $A \subseteq B \implies \bigcap A \supseteq \bigcap B$;
      \item $A \subseteq B \implies A \cap C \subseteq B \cap C$;
      \item $A = B \implies \begin{cases}
               \bigcap A = \bigcap B \\
               A \cap C = B \cap C
            \end{cases}$; 
      \item $(A \cap B) \cap C = A \cap (B \cap C) = \bigcap \{A; B; C\}$;
      \item $\mathfrak{P}(A \cap B) = \mathfrak{P}(A) \cap \mathfrak{P}(B)$.
   \end{enumerate} 
\end{multicols}

\solution

\setcounter{subexercise}{1}
\arabic{subexercise}.Từ định nghĩa,
\begin{align*}
   \bigcap \{A\} &= \left\{z \in \bigcup \{A\} \;\middle|\; \forall a \in \{A\}, z\in a\right\} \\
      &= \left\{z \in A \mid \forall a \in \{A\}, z \in a\right\}.
\end{align*}
Để ý rằng, nếu giả sử $z \in A$, để $a \in \{A\}$ thì $a$ phải bằng $A$, và dẫn đến sự hiển nhiên $z \in a$. Do vậy,
\begin{equation*}
   \forall z, \left(z \in A \implies \forall a \in \{A\}, z \in a\right).
\end{equation*}
Ngoài ra, giả sử $\forall a \in \{A\}, z \in a$. Nhắc lại, $A \in \{A\}$, do đó $z \in A$. Điều này có nghĩa
\begin{equation*}
   \forall z, \left(\forall a \in \{A\}, (z \in a) \implies z \in A\right).
\end{equation*}
Dẫn đến $\forall z, \left(\forall a \in \{A\}, (z \in a) \iff z \in A\right).$ và qua đó $\bigcap \{A\} = \{z \in A \mid z \in A\}$. 

Sử dụng ``$(z) \in A$'' là một vị từ theo hạng thức $z$, từ tiên đề tách, tồn tại tập $y$ để $\forall z, (z \in y \iff z \in y \land z \in A)$. Vậy $y$ có thể là tập nào? Tập $A$ là một ví dụ, do $\forall z, (z \in A \iff z \in A \land z \in A)$. Nhưng vì tính duy nhất của tập xác định từ tiên đề tách, dẫn đến
\begin{equation*}
   \{z \in A \mid z \in A\} = A.
\end{equation*}

Theo tính chất bắc cầu, chúng ta có $\bigcap \{A\} = A$, điều phải chứng minh.

\stepcounter{subexercise}
\arabic{subexercise}
\begin{align*}
   \bigcap \varnothing &= \left\{z \in \bigcup \varnothing \;\middle|\; \forall a \in \varnothing, z\in a\right\} \\ 
   &= \left\{z \in \varnothing \;\middle|\; \forall a \in \varnothing, z\in a\right\}.
\end{align*}
Do đó, với mọi $z$, để $z \in \bigcap \varnothing$ thì trước hết, $z \in \varnothing$. Mặc dù vậy, $z \notin \varnothing$ nên 
\begin{equation*}
   \forall z, z \notin \bigcap \varnothing.
\end{equation*}
Vậy, từ tiên đề tập rỗng, $\bigcap \varnothing = \varnothing$.

\stepcounter{subexercise}
\arabic{subexercise}.Nhắc lại, $A \cap B$ là tập hợp được lấy ra từ tiên đề tách với tập ban đầu là $A$ và vị từ là ``$() \in B$''. Do đó, theo hệ quả \refeq{eq:li_thuyet_tap_hop:tien_de:tien_de_tach:he_qua_tap_con} đã có ở trang \pageref{eq:li_thuyet_tap_hop:tien_de:tien_de_tach:he_qua_tap_con}, $A \cap B \subseteq A$. Điều phải chứng minh.

\stepcounter{subexercise}
\arabic{subexercise}.Giả sử $A \cap B = A$. Xét một phần tử $x \in A$ bất kì. Khi đó, $x \in A \cap B$, và theo định nghĩa của giao, dẫn đến $x \in B$. Do đó, $A \subseteq B$.

Giả sử chiều ngược lại, $A \subseteq B$. Với mọi $e$, nếu $e$ nằm trong $A$ thì $E$ cũng trong $B$, cho chúng ta
$$
\forall e, (e \in A \implies e \in A \land e \in B).
$$
Hiển nhiên có $\forall x, (x \in A \land x \in B \implies x \in A)$. Do vậy, $\forall x, (x \in A \iff x \in A \land x \in B)$. Từ tiên đề tách, có thể thấy $A$ là tập được tách ra từ $A$ với vị từ ``$() \in B$'', nhưng đây cũng là định nghĩa của $A \cap B$. Do đó, $A \cap B = A$.

Kết hợp hai chiều, chúng ta có $A \cap B = A \iff A \subseteq B$, điều phải chứng minh.

\stepcounter{subexercise}
\arabic{subexercise}.Từ tiên đề tách và định nghĩa của giao,
\begin{equation*}
    \forall z, \left(z \in A \cap B \iff z \in A \land z \in B\right).
\end{equation*}
Tương tự, chúng ta cũng có
\begin{equation*}
    \forall z, \left(z \in B \cap A \iff z \in B \land z \in A\right).
\end{equation*}
Theo tính chất bắc cầu của phép đẳng giá, $\forall z, (z \in A \cap B \iff z \in B \cap A)$, và từ tiên đề ngoại diên, $A \cap B = B \cap A$. Điều phải chứng minh.

\stepcounter{subexercise}
\arabic{subexercise}.Nhớ lại, $(A; B) = \{\{A\}; \{A; B\}\}$, dẫn đến $\bigcap (A; B) = \{A\} \cap \{A; B\}$.

Nhận thấy rằng với mọi $w$, để $w$ thuộc $\bigcap (A; B)$ thì $w$ phải thuộc cả $\{A\}$ và $\{A; B\}$. Điều này tương đương với $w = A$. Do vậy, $$\forall w, \left(w \in \bigcap (A; B) \iff w = A\right).$$

Như một kết quả của tiên đề cặp và tiên đề ngoại diên, chỉ tồn tại duy nhất một tập $z$ sao cho $\forall w, (w \in z \iff w = A)$. Do vậy, $\bigcap (A; B) = \{A\}$ và chúng ta có điều phải chứng minh.

\stepcounter{subexercise}
\arabic{subexercise}.Giả sử $A \subseteq B$. Xét $z$ là một phần tử thuộc $\bigcap B$. Khi đó, $z \in b$ với mọi $b \in B$. Do $A \subseteq B$, mọi phần tử $x \in A$ cũng là một phần tử của $B$. Dẫn đến, viết lại các mệnh đề để có
\begin{equation*}
    \begin{cases}
        \forall z \in \bigcap B, \forall b, (b \in B \implies z \in b) \\
        \forall x, (x \in A \implies x \in B)
    \end{cases}.
\end{equation*}
Kết hợp các tính chất của lượng từ rỗng, tính chất phân phối của lượng từ phổ quát với phép hội và tính cụ thể hóa,
\begin{equation*}
    \forall z \in \bigcap B, \forall x, (x \in A \implies z \in x)
\end{equation*}
hay $\forall z, \left(z \in \bigcap B \implies \forall x \in A, z \in x\right)$. 

Ngoài ra, định nghĩa cho chúng ta
$\bigcap A = \{z \in \bigcup A \mid \forall x \in A, z \in x\}$, suy ra \begin{equation*}
    \forall z, \left(z \in \bigcap A \iff z \in \bigcup A \land \forall x \in A, z \in x\right).
\end{equation*}
Nhưng, nhận thấy rằng, với mọi $z$, $\forall x \in A, z \in x \implies \exists x \in A, z \in x$\footnote{Một số người có thể sẽ không chấp nhận cách diễn đạt này với phản biện rằng: ``$A$ có thể là tập rỗng, dẫn đến không tồn tại $x \in A$, và do đó, không thể có $x \in A$ để xét $z \in x$.''. Hãy nhớ lại rằng, theo các khái niệm mà tác giả đã đưa ra, đây chỉ là cách viết tắt. Biểu diễn đầy đủ của mệnh đề đó sẽ là $$
\forall x, (x \in A \implies z \in x) \implies \exists x, (x \in A \implies z \in x)
$$ và điều này thì hoàn toàn đúng.}, dẫn đến $\forall x \in A, z \in x \implies z \in \bigcup A$ là đúng với mọi $z$. Và vì vậy,
\begin{equation}
   \forall z, \begin{cases}
      z \in \bigcap A \iff z \in \bigcup A \land \forall x \in A, z \in x \\
      \forall x \in A, z \in x \implies z \in \bigcup A
   \end{cases}.
   \label{eq:li_thuyet_tap_hop:tien_de:tien_de_tach:7}
\end{equation}

\begin{table}[H]
	\centering
	\caption{Bảng giá trị chân lí của $( P \iff Q \land R ) \land ( R \implies Q ) \implies ( P \iff R )$.}
   \label{tab:li_thuyet_tap_hop:tien_de:tien_de_tach:7}
	\begin{tblr}{
		colspec = {|c|c|c|ccccccccccccc|}
	}
		\hline
		$P$ & $Q$ & $R$ & $(P$ & $\iff$ & $Q$ & $\land$ & $R)$ & $\land$ & $(R$ & $\implies$ & $Q)$ & $\implies$ & $(P$ & $\iff$ & $R)$ \\
		\hline[1.5pt]
		Đ & Đ & Đ & Đ & Đ & Đ & Đ & Đ & Đ & Đ & Đ & Đ & \emphcolor{Đ} & Đ & Đ & Đ \\
		Đ & Đ & S & Đ & S & Đ & S & S & S & S & Đ & Đ & \emphcolor{Đ} & Đ & S & S \\
		Đ & S & Đ & Đ & S & S & S & Đ & S & Đ & S & S & \emphcolor{Đ} & Đ & Đ & Đ \\
		Đ & S & S & Đ & S & S & S & S & S & S & Đ & S & \emphcolor{Đ} & Đ & S & S \\
		S & Đ & Đ & S & S & Đ & Đ & Đ & S & Đ & Đ & Đ & \emphcolor{Đ} & S & S & Đ \\
		S & Đ & S & S & Đ & Đ & S & S & Đ & S & Đ & Đ & \emphcolor{Đ} & S & Đ & S \\
		S & S & Đ & S & Đ & S & S & Đ & S & Đ & S & S & \emphcolor{Đ} & S & S & Đ \\
		S & S & S & S & Đ & S & S & S & Đ & S & Đ & S & \emphcolor{Đ} & S & Đ & S \\
		\hline
	\end{tblr}
\end{table}

Bảng giá trị chân lí \ref{tab:li_thuyet_tap_hop:tien_de:tien_de_tach:7} cho phép chúng ta suy ra từ \refeq{eq:li_thuyet_tap_hop:tien_de:tien_de_tach:7}:
\begin{equation*}
   \forall z, \left(z \in \bigcap A \iff \forall x \in A, z\in x\right).
\end{equation*}

Sử dụng kết quả này để thay thế tương đương, có $\forall z, (z \in \bigcap B \implies z \in \bigcap A)$. Qua đó, $\bigcap B \subseteq \bigcap A$. Từ đây, chúng ta có thể dễ dàng nhìn ra được điều phải chứng minh.

\stepcounter{subexercise}
\arabic{subexercise}. Giả sử $A \subseteq B$. Gọi $x \in A \cap C$ bất kì. Chúng ta có đồng thời $x \in A$ và $x \in C$. Từ giả thiết ban đầu, $x \in A$ có được $x \in B$. Và khi chúng ta có cả $x \in B$ và $x \in C$, $x \in B \cap C$.

Tóm lại, $x \in A \cap C \implies x \in B \cap C$ với mọi $x$, cho nên $A \cap C \subseteq B \cap C$. Qua đó, chúng ta có điều phải chứng minh.

\stepcounter{subexercise}
\arabic{subexercise}. Chúng ta đã chứng minh:
\begin{equation*}
   A \subseteq B \implies
   \begin{cases}
      \bigcap A \supseteq \bigcap B \\
      A \cap C \subseteq B \cap C
   \end{cases}.
\end{equation*}
Và từ đó, chúng ta cũng sẽ có
\begin{equation*}
   A \supseteq B \implies
   \begin{cases}
      \bigcap A \subseteq \bigcap B \\
      A \cap C \supseteq B \cap C
   \end{cases}.
\end{equation*}

Có $x = y \iff (x \subseteq y) \land (x \supseteq y)$ đúng với mọi tập $x$ và $y$. Cho nên, chúng ta dễ dàng thấy được
\begin{equation*}
   A = B \implies \begin{cases}
               \bigcap A = \bigcap B \\
               A \cap C = B \cap C
            \end{cases}.
\end{equation*}
Điểu phải chứng minh.
   
\stepcounter{subexercise}
\arabic{subexercise}. Trước hết, chúng ta sẽ chứng minh $(A \cap B) \cap C = A \cap (B \cap C)$. Thật vậy, với mọi $x$, sử dụng liên tục định nghĩa giao hai tập hợp, chúng ta có
\begin{align*}
   x \in (A \cap B) \cap C &\iff \begin{cases}
      x \in A \cap B \\
      x \in C
   \end{cases}; \\
   x \in (A \cap B) \cap C &\iff \begin{cases}
      x \in A \\
      x \in B \\
      x \in C
   \end{cases}; \\
   x \in (A \cap B) \cap C &\iff \begin{cases}
      x \in A \\
      x \in B \cap C
   \end{cases}; \\
   x \in (A \cap B) \cap C &\iff x \in A \cap (B \cap C).
\end{align*}

Tiếp theo, chúng ta sẽ cần phải khẳng định rằng $(A \cap B) \cap C = \bigcap \{A; B; C\}$. Lấy tùy ý $z \in \bigcap\{A; B; C\}$, điều này tương đương với $z \in \bigcup \{A; B; C\}$ và $\forall y \in \{A; B; C\}, z \in y$. Mà, dễ nhận thấy rằng $\forall y \in \{A; B; C\}, z \in y \implies z \in \bigcup \{A; B; C\}$, cho nên
\begin{equation*}
   z \in \bigcap\{A; B; C\} \iff \forall y \in \{A; B; C\}, z \in y.
\end{equation*}

$\{A; B; C\}$ là tập duy nhất thỏa mãn với mọi $y$, $\left[\begin{array}{l}
   y = A \\
   y = B \\
   y = C   
\end{array}\right.$. Do vậy, với mọi $z$,
\begin{align*}
   \forall y \in \{A; B; C\}, z \in y &\iff \forall y, \left(\left[\begin{array}{l}
   y = A \\
   y = B \\
   y = C   
\end{array} \implies z\in y\right.\right); \\
\forall y \in \{A; B; C\}, z \in y &\iff \forall y, \begin{cases}
   y = A \implies z \in y \\
   y = B \implies z \in y \\
   y = C \implies z \in y 
\end{cases} \equationexplanation{quy tắc chia trường hợp}; \\
\forall y \in \{A; B; C\}, z \in y &\iff \forall y, \begin{cases}
   z \in A \\
   z \in B \\
   z \in C
\end{cases}.
\end{align*}
Bỏ lượng từ rỗng $\forall y$ đi để được $\forall y \in \{A; B; C\}, z \in y \iff \begin{cases}
   z \in A \\
   z \in B \\
   z \in C
\end{cases}$. Đã có $z \in (A \cap B) \cap C \iff \begin{cases}
   z \in A \\
   z \in B \\
   z \in C
\end{cases}$, cho nên bắc cầu để có $$\forall z, z \in \bigcap\{A; B; C\} \iff z \in (A \cap B) \cap C.$$
Do vậy $\bigcap\{A; B; C\} = (A \cap B) \cap C$. Từ đây, theo tính chất bắc cầu của \coloringCrimson{=}, dễ thấy điều phải chứng minh.

\stepcounter{subexercise}
\arabic{subexercise}. Vọi mọi $x$,
\begin{align*}
   x \in \mathfrak{P}(A \cap B) &\iff x \subseteq A \cap B \equationexplanation{định nghĩa tập lũy thừa}; \\
   x \in \mathfrak{P}(A \cap B) &\iff \forall y, (y \in x \implies y \in A \cap B) \equationexplanation{định nghĩa tập con}; \\
   x \in \mathfrak{P}(A \cap B) &\iff \forall y, (y \in x \implies y \in A \land y \in B) \equationexplanation{định nghĩa giao hai tập hợp}; \\
   x \in \mathfrak{P}(A \cap B) &\iff \forall y, \begin{cases}
      y \in x \implies y \in A \\
      y \in x \implies y \in B
   \end{cases} \equationexplanation{$(P \implies (Q \land R)) \iff \begin{cases}
      P \implies Q \\
      P \implies R
   \end{cases}$}; \\
   x \in \mathfrak{P}(A \cap B) &\iff \begin{cases}
      \forall y, y \in x \implies y \in A \\
      \forall y, y \in x \implies y \in B
   \end{cases} \equationexplanation{\begin{tabular}{l}
      tính chất phân phối của lượng từ \\
      phổ quát với phép hội
   \end{tabular}}; \\
   x \in \mathfrak{P}(A \cap B) &\iff x \subseteq A \land x \subseteq B \equationexplanation{định nghĩa tập con}; \\
   x \in \mathfrak{P}(A \cap B) &\iff x \in \mathfrak{P}(A) \land x \in \mathfrak{P}(B) \equationexplanation{định nghĩa tập lũy thừa}; \\
   x \in \mathfrak{P}(A \cap B) &\iff x \in \mathfrak{P}(A) \cap \mathfrak{P}(B) \equationexplanation{định nghĩa giao hai tập hợp}.
\end{align*} 
Và từ đó, $\mathfrak{P}(A \cap B) = \mathfrak{P}(A) \cap \mathfrak{P}(B)$. Điều phải chứng minh.

\exercise Cho $A, B, C, D, E$ là các tập hợp. Chứng minh rằng
\begin{enumerate}
   \item $A \cap (B \cup C) = (A \cap B) \cup (A \cap C)$;
   \item $A \cup (B \cap C) = (A \cup B) \cap (A \cup C)$;
   \item $A \cup B = A \cap C \iff \begin{cases}
      B \subseteq A \\
      A \subseteq C
   \end{cases}$;
   \item $\begin{cases}
      A \cap B = A \cap C \\
      A \cup B = A \cup C
   \end{cases} \iff B = C$;
   \item $A \cap (B \cup (A \cap C)) = A \cap (B \cup C)$;
   \item $A \cup (B \cap (A \cup C)) = A \cup (B \cap C)$;
   \item $A \cap B = C \cap D \implies (A \cup (B \cap C)) \cap (A \cup (B \cap D)) = A$;
   \item $\begin{cases}
      A \cap B = C\cap D\\
      C \cup D = E \\
      C \subseteq A \\
      D \subseteq B \\
      A \subseteq E \\
      B \subseteq E
   \end{cases} \implies \begin{cases}
      C = A\\
      D = B
   \end{cases}$.
\end{enumerate}
\solution

\setcounter{subexercise}{1}
\arabic{subexercise}. Từ những kết quả đã có, với mọi $x$, thực hiện biến đổi như sau
\begin{align*}
   x \in A \cap (B \cup C) &\iff x \in A \land x \in B \cup C; \\
   x \in A \cap (B \cup C) &\iff x \in A \land (x \in B \lor x \in C); \\
   x \in A \cap (B \cup C) &\iff (x \in A \land x \in B) \lor (x \in A \land x \in C) \equationexplanation{tính chất phân phối}; \\
   x \in A \cap (B \cup C) &\iff x \in A \cap B \lor x \in A \cap C; \\
   x \in A \cap (B \cup C) &\iff x \in (A \cap B) \cup (A \cap C).
\end{align*}
Từ đó, dễ có điều phải chứng minh.

\stepcounter{subexercise}
\arabic{subexercise}. Tương tự, với mọi $x$,
\begin{align*}
   x \in A \cup (B \cap C) &\iff x \in A \lor x \in B \cap C; \\
   x \in A \cup (B \cap C) &\iff x \in A \lor (x \in B \land x \in C); \\
   x \in A \cup (B \cap C) &\iff (x \in A \lor x \in B) \land (x \in A \lor x \in C); \\
   x \in A \cup (B \cap C) &\iff x \in A \cup B \land x \in A \cup C; \\
   x \in A \cup (B \cap C) &\iff x \in (A \cup B) \cap (A \cup C)
\end{align*}
và điều phải chứng minh là hiển nhiên.

\stepcounter{subexercise}
\arabic{subexercise}. Đặt $S = A \cup B = A \cap C$. Gọi $x$ là một phần tử thuộc $S$. Theo cách đặt, $x$ cũng thuộc $A \cap C$, dẫn đến $x \in A$. Qua đó, chúng ta được
$$\forall x, x \in S \implies x \in A.$$

Giả sử $x \in A$. Dễ thấy $x \in A \cup B$ cũng đúng, và từ đó, $x \in S$. Do vậy, chúng ta cũng có $$\forall x, x \in A \implies x \in S.$$

Suy ra, $\forall x, x \in A \iff x \in S$. Vậy, $S = A$, cho chúng ta $A \cup B = A \cap C \iff \begin {cases}A = A \cup B \\ A = A \cap C\end{cases}$. Dùng những bài tập trước như hệ quả,
\begin{equation*}
   \begin{cases}
      A = A \cup B \iff B \subseteq A \\
      A = A \cap C \iff A \subseteq C
   \end{cases}.
\end{equation*}
Và vậy, $A \cup B = A \cap C \iff \begin{cases}
   B \subseteq A \\
   A \subseteq C
\end{cases}$. Điều phải chứng minh.

\stepcounter{subexercise}
\arabic{subexercise}. Từ $B = C$ dễ dàng suy ra $\begin{cases}
   A \cap B = A \cap C \\
   A \cup B = A \cup C
\end{cases}$. Chứng minh chiều ngược lại, giả sử $A \cap B = A \cap C$ và $A \cup B = A \cup C$. Khi này,
\begin{align*}
   B  &= (A \cap B) \cup B \\
      &= (A \cap C) \cup B \equationexplanation{giả thiết có $A \cap B = A \cap C$}; \\
      &= (A \cup B) \cap (C \cup B); \\
      &= (C \cup A) \cap (C \cup B) \equationexplanation{$A \cup B = A \cup C = C \cup A$}; \\
      &= C \cup (A \cap B); \\
      &= C \cup (A \cap C) \equationexplanation{$A \cap B = A\cap C$}; \\
      &= C.
\end{align*}
Qua đó, dễ thấy điều phải chứng minh.

\stepcounter{subexercise}
\arabic{subexercise}. Có $A \subseteq B \cup A$ nên $A \cap (B \cup A) = A$. Thực hiện biến đổi để có
\begin{align*}
   A \cap (B \cup (A \cap C)) &= A \cap ((B \cup A) \cap (B \cup C)) \\
   &= (A \cap (B \cup A)) \cap (B \cup C) \\
   &= A \cap (B \cup C).
\end{align*}
Điều phải chứng minh.

\stepcounter{subexercise}
\arabic{subexercise}. Làm tương tự, với để ý rằng $A \supseteq B \cap A$ nên $A \cup (B \cap A) = A$,
\begin{align*}
   A \cup (B \cap (A \cup C)) &= A \cup ((B \cap A) \cup (B \cap C)) \\
   &= (A \cup (B \cap A)) \cup (B \cap C) \\
   &= A \cup (B \cap C).
\end{align*}
Điều phải chứng minh.

\stepcounter{subexercise}
\arabic{subexercise}. Giả sử $A \cap B = C \cap D$,
\begin{align*}
   (A \cup (B \cap C)) \cap (A \cup (B \cap D)) &= A \cup ((B \cap C) \cap (B \cap D)) \\
   &= A \cup (((B \cap C) \cap B) \cap D) \\
   &= A \cup ((B \cap C) \cap D) \equationexplanation{$B \cap C \subseteq B$ nên $(B \cap C) \cap B = B \cap C$} \\
   &= A \cup (B \cap (C \cap D)) \\
   &= A \cup (B \cap (A \cap B)) \equationexplanation{giả thiết} \\
   &= A \cup (A \cap B) \equationexplanation{$A \cap B \subseteq B$ nên $B \cap (A \cap B) = A \cap B$} \\
   &= A.
\end{align*}
Điều phải chứng minh.

\stepcounter{subexercise}
\arabic{subexercise}. Lấy tùy ý $y \in A$. Do $A \subseteq E$ nên $y \in E$, hay $y \in C \cup D$. Suy ra, $y \in C$ hoặc $y \in D$.

Giả sử $y \notin C$, dẫn đến $y \in D$. Kết hợp $D \subseteq B$ để có $y \in B$. $y$ đồng thời thuộc $A$ và $B$ cho khẳng định $y \in A \cap B$. Đề bài cho $A \cap B = C\cap D$ nên $y \in C \cap D$. Vậy $y \in C$, nhưng điều này mâu thuẫn với giả sử ban đầu.

Sử dụng phương pháp chứng minh bằng phản chứng, cần phải có $y\in A \implies y \in C$ với mọi $y$, và đương nhiên, khi đó, $A \subseteq C$. Đã có $C \subseteq A$ nên $C = A$.

Với lập luận tương đương, chúng ta cũng có thể suy ra $D = B$. Kết hợp lại để có được điều phải chứng minh.

\exercise Cho $A, B, C$ là các tập hợp. Chứng minh rằng
\begin{multicols}{2}
   \begin{enumerate}
      \item $A \setminus B \subseteq A$;
      \item $A \setminus B = A \iff A \cap B = \varnothing$;
      \item $A \setminus B = \varnothing \iff A \subseteq B$;
      \item $A \setminus B = A \setminus (A \cap B)$;
      \item $A = (A \setminus B) \cup (A \cap B)$;
      \item $A \subseteq B \iff \begin{cases}
         A \setminus C \subseteq B \setminus C \\
         C \setminus A \supseteq C \setminus B
      \end{cases}$;
      \item $A = B \iff \begin{cases}
         A \setminus C = B \setminus C \\
         C \setminus A = C \setminus B
      \end{cases}$;
      \item $A \setminus (B\cap C) = (A \setminus B) \cup (A \setminus C)$;
      \item $A \setminus (B \cup C) = (A \setminus B) \cap (A \setminus C) = (A \setminus B) \setminus C$;
      \item $A \setminus (B \setminus C) = (A \setminus B) \cup (A \cap C)$;
      \item $(A \setminus B) \setminus (A \setminus C) = (A \setminus B) \cap C = (A \cap C) \setminus B$;
      \item $A \subseteq B \iff \begin{cases}
         A \cap C \subseteq B \cap C \\
         A \setminus C \subseteq B \setminus C
      \end{cases}$.
   \end{enumerate}
\end{multicols}

\solution

\setcounter{subexercise}{1}
\arabic{subexercise}. Có $A \setminus B = \{x \in A \mid x \notin B\}$. Theo hệ quả của tiên đề tách \refeq{eq:li_thuyet_tap_hop:tien_de:tien_de_tach:he_qua_tap_con}, $A \setminus B \subseteq A$. Điều phải chứng minh.

\stepcounter{subexercise}
\arabic{subexercise}. Đầu tiên, giả sử $A \setminus B = A$. Xét một phần tử $x$.
\begin{itemize}
   \item Nếu $x \notin A$, hiển nhiên $x \notin A \cap B$;
   \item Nếu $x \in A$, do $A = A \setminus B$ nên $x \notin B$. Qua đó, $x \notin A \cap B$.
\end{itemize}
Do $x$ chỉ có thể thuộc hoặc không thuộc $A$, cho nên chúng ta có thể suy ra rằng $\forall x, x \notin A \cap B$. Định nghĩa của tập rỗng cũng có biểu thức tương tự. Chỉ có một tập rỗng nên $A \cap B = \varnothing$.

Chứng minh chiều ngược lại, cho $A \cap B = \varnothing$. Gọi $y$ là một phần tử bất kì thuộc $A$. Do $y \notin A \cap B$ (tính chất tập rỗng) nên không thể xảy ra đồng thời $y \in A$ và $y \in B$. Mà $y \in A$ theo giả thiết, dẫn đến $y \notin B$. Qua đó, $y \in A \setminus B$ với mọi $y$ trong $A$, suy ra $A \subseteq A\setminus B$.

Từ phần trước của bài tập này, chúng ta có $A\setminus B \subseteq A$, nên phải có $A = A \setminus B$. 

Kết hợp hai chiều, chúng ta có điều phải chứng minh.

\stepcounter{subexercise}
\arabic{subexercise}. Giả sử $A \setminus B = \varnothing$. Lấy $x \in A$ bất kì. Nếu $x \notin B$, sẽ dẫn đến $x \in A \setminus B$. Nhưng, theo tiên đề tập rỗng, do $A \setminus B = \varnothing$ nên $\forall x, x \notin A \setminus B$, mâu thuẫn. Do vậy, $x \in B$. Chúng ta đã chứng minh, $$\forall x \in A, x \in B$$ hay $A \subseteq B$ trong trường hợp này.

Trường hợp còn lại, giả sử $A \subseteq B$. Xét một phần tử $y$ bất kì.
\begin{itemize}
   \item Nếu $y \notin A$, hiển nhiên $y \notin A \setminus B$;
   \item Nếu $y \in A$, dễ thấy $y$ cũng phải thuộc $B$. Qua đó, $y \notin A \setminus B$.
\end{itemize}
Vậy, với mọi $y$, $y$ luôn không ở trong $A \setminus B$. Hay, $A \setminus B = \varnothing$.

Kết hợp hai chiều, điều phải chứng minh.

\stepcounter{subexercise}
\arabic{subexercise}. Theo định nghĩa của phép hiệu và phép giao:
\begin{align*}
x \in A \setminus (A \cap B) &\iff (x \in A) \land \overline{x \in A \cap B}; \\
x \in A \setminus (A \cap B) &\iff (x \in A) \land \overline{x \in A \land x \in B}; \\
x \in A \setminus (A \cap B) &\iff (x \in A) \land (x \notin A \lor x \notin B) \equationexplanation{định luật Đờ Moóc-gơn}; \\
x \in A \setminus (A \cap B) &\iff ((x \in A) \land (x \notin A)) \lor ((x \in A) \land (x \notin B)) \equationexplanation{tính chất phân phối}.
\end{align*}

Do một phần tử không thể vừa thuộc vừa không thuộc một tập hợp, cho nên
\begin{align*}
   x \in A \setminus (A \cap B) &\iff x \in A \land x \notin B; \\
   x \in A \setminus (A \cap B) &\iff x \in A \setminus B.
\end{align*}

Vậy, $A \setminus (A \cap B) = A \setminus B$. Điều phải chứng minh.

\stepcounter{subexercise}
\arabic{subexercise}. Lấy một phần tử $x$ bất kì thuộc $A$. Xét hai trường hợp xảy ra đối với tập hợp $B$:
\begin{itemize}
    \item Nếu $x \in B$, vì $x \in A$ và $x \in B$, ta có $x \in A \cap B$.\\
    Suy ra $x \in (A \setminus B) \cup (A \cap B)$.
    
    \item Nếu $x \notin B$, vì $x \in A$ và $x \notin B$, ta có $x \in A \setminus B$.\\
    Suy ra $x \in (A \setminus B) \cup (A \cap B)$.
\end{itemize}
Trong cả hai trường hợp, đều có $x \in (A \setminus B) \cup (A \cap B)$.\\
Do đó,
\[ A \subseteq (A \setminus B) \cup (A \cap B). \]

Lấy một phần tử $y \in (A \setminus B) \cup (A \cap B)$. Theo định nghĩa phép hợp,
\begin{itemize}
   \item $x \in A \setminus B$. Theo định nghĩa phép hiệu, $x \in A$ (và $x \notin B$);  
   \item Hoặc $x \in A \cap B$. Theo định nghĩa phép giao, $x \in A$ (và $x \in B$).
\end{itemize}
Trong cả hai trường hợp, chúng ta đều có $x \in A$.\\
Vì vậy,
\[ (A \setminus B) \cup (A \cap B) \subseteq A. \]

Vậy hai tập hợp là bằng nhau, điều phải chứng minh.

\stepcounter{subexercise}
\arabic{subexercise}. Trước hết, giả sử $A \subseteq B$. Lấy một phần tử $x$ bất kì thuộc vào tập hợp $A \setminus C$. Theo định nghĩa của phép hiệu tập hợp,
$ x \in A \land x \notin C. $
Vì $A \subseteq B$ và $x \in A$, nên
$ x \in B. $
Từ đó, có đồng thời:
\[ x \in B \land x \notin C. \]
Theo định nghĩa của phép hiệu, điều này có nghĩa là $ x \in B \setminus C $
Vì mọi phần tử $x \in A \setminus C$ đều thuộc $B \setminus C$, kết luận được
\[ A \setminus C \subseteq B \setminus C. \]

Lấy một phần tử $y$ bất kì thuộc vào tập hợp $C \setminus B$. Chúng ta có $ y \in C \land y \notin B $. Cần chứng minh rằng $y \notin A$. Giả sử ngược lại rằng $y \in A$. Vì $A \subseteq B$, nếu $y \in A$ thì ta bắt buộc phải có $y \in B$. Điều này mâu thuẫn trực tiếp với điều kiện ban đầu là $y \notin B$.
Do đó, giả thiết phản chứng là sai. Dẫn đến $ y \notin A $. Như vậy, chúng ta có đồng thời
\[ y \in C \land y \notin A. \]
Theo định nghĩa của phép hiệu, điều này có nghĩa là $ y \in C \setminus A $. Vậy
\[ C \setminus B \subseteq C \setminus A \quad \text{hay} \quad C \setminus A \supseteq C \setminus B. \]

Kết hợp hai kết quả, chúng ta có 
\begin{equation}
   A \subseteq B \implies \begin{cases}
      A \setminus C \subseteq B \setminus C \\
      C \setminus A \supseteq C \setminus B
   \end{cases}.
   \label{eq:li_thuyet_tap_hop:tien_de:tien_de_tach:so_sanh_hieu.1}
\end{equation}

Sau đó, chúng ta giả sử chiều ngược lại rằng $\begin{cases}
   A \setminus C \subseteq B \setminus C \\
   C \setminus A \supseteq C \setminus B
\end{cases}$. Sử dụng chứng minh bằng phương pháp phản chứng, giả sử $A \not \subseteq B$, tương đương với $\overline{\forall z, (z \in A \implies z \in B)}$. Theo lập luận lô-gích thông thường, $\exists z, (z \in A \land z \notin B)$. Cụ thể hóa phần tử đó và gọi nó là $w$. Xét hai trường hợp:
\begin{itemize}
   \item Trường hợp \coloringCrimson{$w \in C$}: Khi này, chúng ta có $w \in C \land w \in A$ và $w \in C \land w \notin B$. Dẫn đến
   \begin{equation*}
      \begin{cases}
         w \notin C \setminus A\\
         w \in C \setminus B
      \end{cases}.
   \end{equation*} Chúng ta thấy tồn tại phần tử $w$ trong $C \setminus B$ không thuộc $C \setminus A$, nhưng giả thiết $C \setminus A \supseteq C \setminus B$ lại mâu thuẫn với điều này.
   \item Trường hợp \coloringDodgerblue{$w \notin C$}: Chúng ta cũng dễ dàng suy ra được $w \in A \setminus C$ và $w \notin B \setminus C$, nhưng $A \setminus C \subseteq B \setminus C$ cũng tạo ra mâu thuẫn.
\end{itemize}
Do mọi trường hợp đều ra mâu thuẫn, nên điều giả sử ban đầu là sai. Do đó, $A \subseteq B$. Vậy,
\begin{equation}
   \begin{cases}
      A \setminus C \subseteq B \setminus C \\
      C \setminus A \supseteq C \setminus B
   \end{cases} \implies A \subseteq B.
   \label{eq:li_thuyet_tap_hop:tien_de:tien_de_tach:so_sanh_hieu.2}
\end{equation}

Từ \refeq{eq:li_thuyet_tap_hop:tien_de:tien_de_tach:so_sanh_hieu.1} và \refeq{eq:li_thuyet_tap_hop:tien_de:tien_de_tach:so_sanh_hieu.2}, chúng ta có điều phải chứng minh.

\stepcounter{subexercise}
\arabic{subexercise}. Phần trước cho 
\begin{equation*}
   A \subseteq B \iff \begin{cases}
      A \setminus C \subseteq B \setminus C \\
      C \setminus A \supseteq C \setminus B
   \end{cases}.
\end{equation*} Sử dụng chính kết quả này, chúng ta cũng có
\begin{equation*}
   B \subseteq A \iff \begin{cases}
      B \setminus C \subseteq A \setminus C \\
      C \setminus B \subseteq C \setminus A
   \end{cases}.
\end{equation*}
Qua đó,
\begin{equation*}
   A \subseteq B \land B \subseteq A \iff \begin{cases}
      A \setminus C \subseteq B \setminus C \\
      C \setminus A \supseteq C \setminus B \\
      B \setminus C \subseteq A \setminus C \\
      C \setminus B \subseteq C \setminus A
   \end{cases}
\end{equation*}
và dễ dàng suy ra được
\begin{equation*}
   A = B \iff \begin{cases}
      A \setminus C = B \setminus C \\
      C \setminus A = C \setminus B
   \end{cases}.
\end{equation*}
Điều phải chứng minh.

\stepcounter{subexercise}
\arabic{subexercise}. Lấy $x \in A \setminus (B \cap C)$. Khi này,
\begin{align*}
   x \in A \setminus (B \cap C) &\iff x \in A \land x \notin B \cap C; \\
   x \in A \setminus (B \cap C) &\iff x \in A \land \overline{x \in B \cap C}; \\
   x \in A \setminus (B \cap C) &\iff x \in A \land \overline{x \in B \land x \in C}; \\
   x \in A \setminus (B \cap C) &\iff x \in A \land (x \notin B \lor x \notin C) \equationexplanation{định luật Đờ Moóc-gan}; \\
   x \in A \setminus (B \cap C) &\iff (x \in A \land x \notin B) \lor (x \in A \land x\notin C) \equationexplanation{tính chất phân phối}; \\
   x \in A \setminus (B \cap C) &\iff x \in A \setminus B \lor x \in A \setminus C; \\
   x \in A \setminus (B \cap C) &\iff x \in (A \setminus B) \cup (A \setminus C).
\end{align*}
Từ đó, dễ thấy điều phải chứng minh.

\stepcounter{subexercise}
\arabic{subexercise}. Đầu tiên, cần phải chứng minh $A \setminus (B \cup C) = (A \setminus B) \cap (A \setminus C)$. Với lập luận tương tự như phần trước, với mọi $x$,
\begin{align*}
   x \in A \setminus (B \cup C) &\iff x \in A \land x \notin B \cup C; \\
   x \in A \setminus (B \cup C) &\iff x \in A \land \overline{x \in B \cup C}; \\
   x \in A \setminus (B \cup C) &\iff x \in A \land \overline{x \in B \lor x \in C}; \\
   x \in A \setminus (B \cup C) &\iff x \in A \land (x \notin B \land x \notin C) \equationexplanation{định luật Đờ Moóc-gan}; \\
   x \in A \setminus (B \cup C) &\iff (x \in A \land x \notin B) \land (x \in A \land x\notin C); \\
   x \in A \setminus (B \cup C) &\iff x \in A \setminus B \land x \in A \setminus C; \\
   x \in A \setminus (B \cup C) &\iff x \in (A \setminus B) \cap (A \setminus C)
\end{align*}
sẽ dẫn đến $A \setminus (B \cup C) = (A \setminus B) \cap (A \setminus C)$ như chúng ta cần.

Chúng ta cũng có thể lập luận theo một hướng khác:
\begin{align*}
   x \in A \setminus (B \cup C) &\iff x \in A \land (x \notin B \land x \notin C); \\
   x \in A \setminus (B \cup C) &\iff (x \in A \land x \notin B) \land x \notin C; \\
   x \in A \setminus (B \cup C) &\iff x \in A \setminus B \land x \notin C; \\
   x \in A \setminus (B \cup C) &\iff x \in (A \setminus B) \setminus C.
\end{align*}
Và từ đó, chúng ta cũng có $A \setminus (B \cup C) = (A \setminus B) \setminus C$. Bắc cầu hai kết quả để có điều phải chứng minh.

\stepcounter{subexercise}
\arabic{subexercise}. Với mọi $x$,
\begin{align*}
   x \in A \setminus (B \setminus C) &\iff x \in A \land x \notin B \setminus C; \\
   x \in A \setminus (B \setminus C) &\iff x \in A \land \overline{x \in B \land x \notin C}; \\
   x \in A \setminus (B \setminus C) &\iff x \in A \land \left(x \notin B \lor x \in C\right); \\
   x \in A \setminus (B \setminus C) &\iff (x \in A \land x \notin B) \lor (x \in A \land x \in C); \\
   x \in A \setminus (B \setminus C) &\iff x \in A \setminus B \lor x \in A \cap C; \\
   x \in A \setminus (B \setminus C) &\iff x \in (A \setminus B) \cup (A \cap C).
\end{align*}
Do đó, $A \setminus (B \setminus C) = (A \setminus B) \cup (A \cap C)$. Điều phải chứng minh.

\stepcounter{subexercise}
\arabic{subexercise}. Với mọi $x$,
\begin{align*}
   x \in (A \setminus B) \setminus (A \setminus C) &\iff x \in A \setminus B \land x \notin A \setminus C; \\
   x \in (A \setminus B) \setminus (A \setminus C) &\iff (x \in A \land x \notin B) \land \overline{x \in A \land x \notin C}; \\
   x \in (A \setminus B) \setminus (A \setminus C) &\iff (x \in A \land x \notin B) \land \overline{x \in A \land x \notin C}; \\
   x \in (A \setminus B) \setminus (A \setminus C) &\iff x \in A \land x \notin B \land (x \notin A \lor x \in C); \\
   x \in (A \setminus B) \setminus (A \setminus C) &\iff (x \in A \land x \notin B \land x \notin A) \lor (x \in A \land x \notin B \land x \in C).
\end{align*}
Nhận thấy rằng $x \in A \land x \notin A$ là luôn sai để không có mâu thuẫn, cho nên 
\begin{align*}
   x \in (A \setminus B) \setminus (A \setminus C) &\iff x \in A \land x \notin B \land x \in C; \\
   x \in (A \setminus B) \setminus (A \setminus C) &\iff \begin{cases}
      (x \in A \land x \notin B) \land x \in C \\
      (x \in A \land x \in C) \land x \notin B
   \end{cases}; \\
   x \in (A \setminus B) \setminus (A \setminus C) &\iff \begin{cases}
      x \in A \setminus B \land x \in C \\
      x \in A \cap C \land x \notin B
   \end{cases}; \\
   x \in (A \setminus B) \setminus (A \setminus C) &\iff \begin{cases}
      x \in (A \setminus B) \cap C \\
      x \in (A \cap C) \setminus B
   \end{cases}.
\end{align*}
Từ đó, điều phải chứng minh là khá hiển nhiên.

\stepcounter{subexercise}
\arabic{subexercise}. Chứng minh chiều thuận từ $A \subseteq B$ sang $\begin{cases}
   A \cap C \subseteq B \cap C \\
   A \setminus C \subseteq B \setminus C
\end{cases}$ đã được tác giả trình bày ở những bài tập trước. Chứng minh chiều ngược lại, chúng ta giả thiết đã có $\begin{cases}
   A \cap C \subseteq B \cap C \\
   A \setminus C \subseteq B \setminus C
\end{cases}$. Sử dụng chứng minh bằng phương pháp phản chứng, giả sử $A \not \subseteq B$, tương đương với $\exists z, (z \in A \land z \notin B)$. Gọi phần tử thỏa mãn là $w$. Xét hai trường hợp:
\begin{itemize}
   \item Trường hợp \coloringCrimson{$w \in C$}: Khi này, chúng ta có $w \in A \land w \in C$ và $w \notin B \land w \in C$. Dẫn đến
   \begin{equation*}
      \begin{cases}
         w \in A \cap C\\
         w \notin B \cap C
      \end{cases},
   \end{equation*} nhưng giả thiết $A \cap C \supseteq B \cap C$ lại mâu thuẫn với điều này.
   \item Trường hợp \coloringDodgerblue{$w \notin C$}: Chúng ta cũng dễ dàng suy ra được $w \in A \setminus C$ và $w \notin B \setminus C$, nhưng $A \setminus C \subseteq B \setminus C$ cũng tạo ra mâu thuẫn.
\end{itemize}
Do mọi trường hợp đều ra mâu thuẫn, nên điều giả sử ban đầu là sai. Do đó, cần phải có $A \subseteq B$.

Kết hợp hai chiều, chúng ta có điều phải chứng minh.

\exercise Cho $A, B, C$ là các tập hợp. Chứng minh rằng
\begin{multicols}{2}
   \begin{enumerate}
      \item $A \vartriangle B = (A \cup B) \setminus (A \cap B)$;
      \item $A \vartriangle B = B \vartriangle A$;
      \item $A \vartriangle B = A \iff B = \varnothing$;
      \item $A \vartriangle B = \varnothing \iff B = A$;
      \item $A \vartriangle B = C \implies \begin{cases}
         B \vartriangle C = A \\
         C \vartriangle A = B
      \end{cases}$;
      \item $A = B \implies A \vartriangle C = B \vartriangle C$;
      \item $(A \vartriangle B) \vartriangle C = A \vartriangle (B \vartriangle C)$.
   \end{enumerate}
\end{multicols}

\solution

\setcounter{subexercise}{1}
\arabic{subexercise}. Gọi $x \in A \vartriangle B$. Theo định nghĩa, $x \in (A \setminus B) \cup (B \setminus A)$, hay $x \in A \setminus B$ hoặc $x \in B \setminus A$. Qua đó, $x$ cần phải thuộc $A$ hoặc $B$. Do đó $x \in A\cup B$.

Nếu $x \in A\cap B$, điều này có nghĩa là đồng thời $x \in A$ và $x \in B$, dẫn đến $x \in A \setminus B$ và $x \in B \setminus A$ đều không thể xảy ra, mâu thuẫn. Do vậy, $x \notin A \cap B$. Và từ đó, $x \in (A \cup B) \setminus (A \cap B)$ là điều hiển nhiên. Vậy,
\begin{equation*}
   A \vartriangle B \subseteq (A \cup B) \setminus (A \cap B).
\end{equation*}

Ngược lại, với mọi $y$, nếu $y \in (A \cup B) \setminus (A \cap B)$. Chúng ta có $y \in A \cup B$.
\begin{itemize}
   \item Nếu $y \in A$, thấy được rằng $y \notin A \cap B$ nên cần phải có phủ định của $y \in A \land y \in B$. Điều này tương đương với $y \notin A \lor y \notin B$. Đang xét $y \in A$ nên $y \notin B$, dẫn ra được $y \in A \setminus B$.
   \item Nếu $y \in B$, với lập luận tương tự, $y \in B \setminus A$.
\end{itemize}
Do đó, $y \in (A \setminus B) \cup (B \setminus A)$. Vậy,
\begin{equation*}
   (A \cup B) \setminus (A \cap B) \subseteq A \vartriangle B.
\end{equation*}

Từ hai hướng, kết hợp chúng lại để có điều phải chứng minh.

\stepcounter{subexercise}
\arabic{subexercise}. Điều phải chứng minh là khá hiển nhiên, do 
\begin{align*}
   A \vartriangle B &= (A \setminus B) \cup (B \setminus A) \\
      &= (B \setminus A) \cup (A \setminus B) \\
      &= B \vartriangle A.
\end{align*}

